% !TeX root = ../../Book5.tex
% 纲要标题：## 第5章　切空间、Kähler 微分与局部方程
% 纲要位置：几何学计划纲要.md，第 5185 行。
% 正文状态：初稿；尚未正式审核。

% 纲要内容开始
% 
% > 本章目标与证明范围：把切向量、微分模、切锥及形式方程放在不同层次。正则性的定义在此出现，完整局部结构在后面证明；光滑性留待相对态射理论。
% 
% 纲要内容结束

\chapter{切空间、Kähler 微分与局部方程}
\label{g5:ch:05}
\BookChapterNumber{g5:ch:05}

\begin{chapterabstract}
	本章用导子和微分模统一描述代数方程的一阶变化。普遍导子把求导规则转化为模的泛性质，余法正合列和传递正合列分别处理商环与复合映射。在有理点处，微分模的纤维、极大理想模平方的对偶及双数点给出同一个切空间，Jacobian 矩阵使它成为可计算的线性方程组。正则性比较嵌入维数与局部维数；不可分扩域说明它与相对于基域的性质有所不同。伴随分次、切锥和形式方程逐层恢复被线性化舍去的信息。
\end{chapterabstract}

\begin{learninggoals}
	\begin{itemize}
		\item 从导子的加法与乘法规则构造微分模，并计算多项式环、局部化及简单域扩张的微分。
		\item 证明余法正合列、传递正合列与微分模的基变换公式，识别左端核的来源。
		\item 在有理点处建立双数点、导子和极大理想模平方之间的对应，用 Jacobian 矩阵计算切空间。
		\item 依据局部维数判定具体局部环的正则性，计算不可分扩域带来的变化。
		\item 求出局部方程的最低次初始理想和切锥，区分一阶信息、有限阶截断与相容的全部形式邻域。
	\end{itemize}
\end{learninggoals}

在原点处，尖点 \(y^2=x^3\) 与两条相交直线 \(xy=0\) 的方程都没有一次项。仅保留一次项时，两者都给出整个二维切空间；保留最低次非零项时，前者给出 \(y^2=0\)，后者给出 \(xy=0\)。一个是带有加厚结构的重直线，一个是两条不同的直线。线性化提供共同的容纳空间，而最低次方程开始区别它们的局部形状。

\ProseParagraphBreak

求导规则把这个观察写成代数。若坐标只发生满足 \(\varepsilon^2=0\) 的小变化 \(x_i\mapsto a_i+v_i\varepsilon\)，方程变为
\[
	f(a+v\varepsilon)
	=f(a)+\varepsilon\sum_i\frac{\partial f}{\partial x_i}(a)v_i.
\]
导子记录其中 \(\varepsilon\) 的系数，微分模收集全部这样的线性测量。这一构造能直接处理带幂零元的环，因而也能把加厚点纳入讨论。随着 \(\varepsilon^2=0\) 改为更高次的截断条件，二次及更高次方程重新出现，切向量能否继续延拓便成为新的问题。

\ProseParagraphBreak

本章使用第二章的模、张量积、局部化与 Nakayama 引理，以及第三章的局部维数和嵌入维数。第四章的正规化在曲线例子中提供对照：正规化补入函数，微分模和切锥则检测已有函数之间的局部关系。微分模及两个基本正合列可对照 \cite[第25节，第190—196页]{Matsumura1989}。正则局部环的结构定理在第九章展开；相对光滑性在第十五章建立，本章先计算说明这一区别所需的不可分例子。

% !TeX root = ../../Book5.tex
% 纲要标题：### 5.1 普遍导子与微分模
% 纲要位置：几何学计划纲要.md，第 5189 行。
% 正文状态：初稿；尚未正式审核。

% 纲要内容开始
% 
% * 导子的基环线性
% * Kähler 微分的构造和泛性质
% * 多项式环、局部化与张量积的微分
% * 可分与不可分扩张的微分
% 
% 纲要内容结束

\section{普遍导子与微分模}
\label{g5:sec:0501}

多项式的求导法则只使用加法、乘法和整数系数。例如 \(\mathrm{d}(xy)=x\,\mathrm{d}y+y\,\mathrm{d}x\)，并不需要先给 \(x,y\) 指定实数值。这使我们能够在任意交换环上讨论微分。不过，哪些元素应当看作常数，必须由基环指定；同一个环相对于不同基环可以有完全不同的微分。我们先把求导规则抽象为取值于模的映射，再构造一个模，把所有这样的映射同时表达出来。

\subsection{导子的基环线性}
\label{g5:subsec:0501:01}

设 \(A\to B\) 是交换含幺环的同态。这里的 \(A\) 不必是域，同态也不必单射；说一个 \(A\) 元素在 \(B\) 中为常数，是指它在这个同态下的像。微分的值可以属于 \(B\) 本身，也可以属于另一个 \(B\) 模 \(M\)。后一个选择容许我们记录只在某个点、某个商模或某个局部化中可见的变化。

\begin{definition}\label{g5:def:derivation}
	设 \(\varphi:A\to B\) 为交换含幺环同态，\(M\) 为含幺 \(B\) 模。加法映射 \(D:B\to M\) 若满足
	\[
		D(bb')=bD(b')+b'D(b),\qquad
		D(\varphi(a))=0
		\quad(b,b'\in B,\ a\in A),
	\]
	就称为从 \(B\) 到 \(M\) 的一个 \(A\)\term[daozi]{导子}[derivation]。所有这样的导子组成的集合记为 \(\operatorname{Der}_A(B,M)\)。
\end{definition}
\symidx{\(\operatorname{Der}_A(B,M)\)：从 \(B\) 到 \(M\) 的 \(A\) 导子}

导子总满足 \(D(1)=0\)：对 \(1=1\cdot1\) 使用乘法法则，得到 \(D(1)=2D(1)\)，再在加法群中消去一项即可，不需要反演 \(2\)。因此 \(D(n\cdot1)=0\) 对所有整数 \(n\) 成立。进一步，由 \(D(\varphi(a))=0\) 得
\[
	D(\varphi(a)b)=\varphi(a)D(b),
\]
所以每个 \(A\) 导子都是 \(A\) 线性的。反过来，一个满足乘法法则的加法映射若为 \(A\) 线性，则 \(D(\varphi(a))=\varphi(a)D(1)=0\)。定义中的常数条件因而恰好表达了基环线性。

\ProseParagraphBreak

这里的线性针对 \(A\)，乘法法则则针对 \(B\)。如果再要求 \(D\) 为 \(B\) 线性，就会有 \(D(b)=bD(1)=0\)，只剩零导子。例如 \(B=A[t]\) 上的形式求导满足 \(D(t)=1\)，所以在 \(A\ne0\) 时并不是 \(B\) 线性映射。另一方面，导子本身可以乘以 \(B\) 中的系数：逐点定义
\[
	(D+E)(b)=D(b)+E(b),\qquad
	(cD)(b)=cD(b).
\]
交换性保证 \(cD\) 仍满足乘法法则，因此 \(\operatorname{Der}_A(B,M)\) 是一个 \(B\) 模。若 \(h:M\to N\) 为 \(B\) 线性映射，则 \(h\circ D\) 仍为 \(A\) 导子。

\begin{example}\label{g5:exm:derivation-coefficients}
	设 \(k\) 为域，\(B=k[t]\)。形式求导
	\[
		\partial_t\!\left(\sum_i a_it^i\right)=\sum_i i a_it^{i-1}
	\]
	是一个取值于 \(B\) 的 \(k\) 导子；每个 \(g(t)\partial_t\) 也是。若固定 \(a\in k\)，把 \(k\) 看作由求值 \(B\to k\)、\(f\mapsto f(a)\) 给出的 \(B\) 模，则
	\[
		D_a(f)=f'(a)
	\]
	是取值于这个模的 \(k\) 导子，因为
	\[
		D_a(fg)=f(a)D_a(g)+g(a)D_a(f).
	\]
	右边的系数作用是求值后的作用，故同一个形式求导运算能够产生不同目标模中的导子。
\end{example}

反复应用乘法法则，对任意 \(n\geq1\) 都有
\[
	D(b^n)=n b^{n-1}D(b).
\]
若 \(b\) 是单位，则对 \(bb^{-1}=1\) 求导得到
\begin{equation}\label{g5:eq:derivative-inverse}
	D(b^{-1})=-b^{-2}D(b).
\end{equation}
这些式子已经预示了局部化与域扩张的计算。在特征 \(p>0\) 的环上，尤其有 \(D(b^p)=0\)。于是导子可能把许多不属于基环的元素也送到零。例如 \(\mathbb F_p[t]\) 上的 \(\partial_t\) 杀死全部 \(t^p\) 的多项式。反过来，逐项比较系数可知
\[
	\ker(\partial_t:\mathbb F_p[t]\to\mathbb F_p[t])
	=\mathbb F_p[t^p].
\]
所以“微分为零”与“属于指定基环”是两个不同的条件。

\subsection{Kähler 微分的构造和泛性质}
\label{g5:subsec:0501:02}

我们希望先形成一个只由 \(A\to B\) 决定的模，然后把任意 \(D:B\to M\) 改写成这个模到 \(M\) 的线性映射。为此给每个 \(b\in B\) 引入一个形式符号 \(\mathrm{d}b\)，并只施加导子必须满足的关系。这样的构造把非线性的乘法法则放进了源模，取值于各个 \(M\) 的问题便转化为普通模同态。

\begin{definition}\label{g5:def:kahler-differentials}
	设 \(A\to B\) 为交换含幺环同态。若一个 \(B\) 模 \(\Omega_{B/A}\) 连同 \(A\) 导子
	\[
		\mathrm{d}_{B/A}:B\longrightarrow\Omega_{B/A}
	\]
	满足：对任意 \(B\) 模 \(M\) 及任意 \(D\in\operatorname{Der}_A(B,M)\)，存在唯一的 \(B\) 线性映射 \(\widetilde D:\Omega_{B/A}\to M\)，使 \(D=\widetilde D\circ\mathrm{d}_{B/A}\)，则称 \(\Omega_{B/A}\) 为 \(B\) 在 \(A\) 上的\term[Kahlerweifenmo]{Kähler 微分模}[module of Kähler differentials]，称 \(\mathrm{d}_{B/A}\) 为\term[pubiandaozi]{普遍导子}[universal derivation]。\footnote{Erich Kähler（1906--2000），德国数学家，研究微分方程、复几何与数学物理，引入了以其姓氏命名的度量结构。}基环明确时，把 \(\mathrm{d}_{B/A}b\) 简写为 \(\mathrm{d}b\)。
\end{definition}
\symidx{\(\Omega_{B/A}\)：\(B\) 在 \(A\) 上的 Kähler 微分模}
\symidx{\(\mathrm{d}_{B/A}\)：相对于基环 \(A\) 的普遍导子}

泛性质由下图表达。斜箭头与横箭头为 \(A\) 导子，虚线箭头为唯一的 \(B\) 线性映射；交换条件要求每个 \(b\) 的像满足 \(\widetilde D(\mathrm{d}b)=D(b)\)。
\[
	\begin{tikzcd}[column sep=large,row sep=large]
		B\arrow[r,"\mathrm{d}_{B/A}"]\arrow[dr,"D"']&
		\Omega_{B/A}\arrow[d,dashed,"\widetilde D"]\\
		&M
	\end{tikzcd}
\]

\begin{proposition}\label{g5:prop:differentials-existence}
	对任意交换含幺环同态 \(A\to B\)，Kähler 微分模及普遍导子存在。两个满足泛性质的模连同导子之间，存在唯一与导子相容的 \(B\) 模同构。对任意 \(B\) 模 \(M\)，有自然的 \(B\) 模同构
	\begin{equation}\label{g5:eq:differentials-represent-derivations}
		\operatorname{Hom}_B(\Omega_{B/A},M)
		\xrightarrow{\ \sim\ }\operatorname{Der}_A(B,M),
		\qquad h\longmapsto h\circ\mathrm{d}_{B/A}.
	\end{equation}
\end{proposition}
\begin{proof}
	取以 \(B\) 的底集合为指标的自由 \(B\) 模
	\[
		F=\bigoplus_{b\in B}B e_b.
	\]
	这里每个元素都是有限个 \(e_b\) 的线性组合。设 \(N\subseteq F\) 为下列元素生成的子模：
	\[
		\begin{gathered}
			e_{b+b'}-e_b-e_{b'},\qquad
			e_{bb'}-b e_{b'}-b'e_b,\\
			e_{\varphi(a)}
			\qquad(b,b'\in B,\ a\in A).
		\end{gathered}
	\]
	令 \(\Omega_{B/A}=F/N\)，定义 \(\mathrm{d}_{B/A}(b)=e_b+N\)。三类关系分别保证加法性、乘法法则及在 \(A\) 上为零，所以这是一个 \(A\) 导子。

	若 \(D:B\to M\) 为 \(A\) 导子，赋值 \(e_b\mapsto D(b)\) 唯一延伸为 \(F\to M\) 的 \(B\) 线性映射。导子的三条性质使它杀死 \(N\)，于是下降为 \(\widetilde D:F/N\to M\)。由于 \(\mathrm{d}b\)、\(b\in B\)，生成整个商模，条件 \(\widetilde D(\mathrm{d}b)=D(b)\) 又唯一决定了这个映射。这就证明存在性与泛性质。

	若 \((P,\delta)\) 也满足同一泛性质，分别以 \(P\) 和 \(\Omega_{B/A}\) 为测试模，得到
	\[
		u:\Omega_{B/A}\to P,\qquad v:P\to\Omega_{B/A},
		\qquad u\circ\mathrm{d}=\delta,\quad v\circ\delta=\mathrm{d}.
	\]
	于是 \(vu\circ\mathrm{d}=\mathrm{d}\)。恒等映射也满足这个等式，由唯一性 \(vu=1\)；同理 \(uv=1\)。相容同构因此存在且唯一。

	最后，\eqref{g5:eq:differentials-represent-derivations}的双射就是泛性质，其与加法及 \(B\) 标量相容可逐点检验。若 \(M\to M'\) 为 \(B\) 线性映射，两侧都通过复合得到相应映射，且上述双射保持复合。这就是这里所说的自然性。
\end{proof}

这个定义并不要求先给 \(B\) 选择有限组生成元，也不要求 \(A\) 或 \(B\) 为 Noether 环。如果 \(A\to B\) 为满射，每个 \(b\in B\) 都来自基环，所以任意 \(A\) 导子均为零，进而 \(\Omega_{B/A}=0\)。特别地，\(\Omega_{A/A}=0\) 以及 \(\Omega_{(\mathbb Z/n\mathbb Z)/\mathbb Z}=0\)。这些结论都直接取决于哪些元素被规定为常数。

\ProseParagraphBreak

环同态也自然地作用在微分上。具体地，给定 \(A\) 代数同态 \(u:B\to C\)，映射 \(b\mapsto\mathrm{d}_{C/A}(u(b))\) 是取值于 \(C\) 模 \(\Omega_{C/A}\) 的 \(A\) 导子。先用泛性质，再用标量扩张，便得到
\begin{equation}\label{g5:eq:differentials-natural-map}
	C\otimes_B\Omega_{B/A}\longrightarrow\Omega_{C/A},
	\qquad c\otimes\mathrm{d}b\longmapsto c\,\mathrm{d}(u(b)).
\end{equation}
对复合的环同态，所得映射的复合在每个 \(\mathrm{d}b\) 上仍给出相同的微分，因此与复合相容。形式符号构造及这一泛性质可对照\cite[第25节，第190--193页]{Matsumura1989}；该书采用乘法映射 \(B\otimes_A B\to B\) 的核模平方来给出同一个普遍对象。

\subsection{多项式环、局部化与张量积的微分}
\label{g5:subsec:0501:03}

多项式环中的变量没有代数关系，因而可以独立指定每个变量的导数。局部化后的分母则由乘法逆元强制决定导数。这两点把许多微分模的计算化为自由模运算与分式运算。

\begin{proposition}\label{g5:prop:polynomial-differentials}
	设 \(A\) 为交换含幺环，\(B=A[x_i\mid i\in I]\) 为一组独立变量的多项式环，指标集 \(I\) 可以无限。则
	\[
		\Omega_{B/A}\cong\bigoplus_{i\in I}B\,\mathrm{d}x_i,
		\qquad
		\mathrm{d}f=\sum_{i\in I}\frac{\partial f}{\partial x_i}\,\mathrm{d}x_i.
	\]
	对每个具体的 \(f\)，右侧只有有限项非零。对任意 \(B\) 模 \(M\)，赋值 \(D\mapsto(D(x_i))_{i\in I}\) 给出
	\[
		\operatorname{Der}_A(B,M)\cong\prod_{i\in I}M.
	\]
\end{proposition}
\begin{proof}
	首先对单项式反复使用乘法法则。对
	\(x_1^{a_1}\cdots x_r^{a_r}\)，得到
	\[
		D(x_1^{a_1}\cdots x_r^{a_r})
		=\sum_{j=1}^r a_jx_1^{a_1}\cdots x_j^{a_j-1}
		\cdots x_r^{a_r}D(x_j),
	\]
	其中 \(a_j=0\) 的项省略。再用加法性与在 \(A\) 上为零，便知每个导子都必须满足
	\[
		D(f)=\sum_i\frac{\partial f}{\partial x_i}D(x_i).
	\]
	所以它由各个 \(D(x_i)\) 唯一决定。

	反过来，任意指定 \(m_i\in M\)，用
	\(D(f)=\sum_i(\partial f/\partial x_i)m_i\) 定义 \(D\)。多项式的每一项只含有限个变量，整个多项式也只有有限项，所以公式始终是有限和。单项式指数相加给出形式偏导的乘法法则，再按加法展开就得到一般的乘法法则；系数在 \(A\) 中时偏导为零。因此 \(D\) 确为 \(A\) 导子。

	令 \(E=\bigoplus_i B\varepsilon_i\)，并取
	\(\delta(f)=\sum_i(\partial f/\partial x_i)\varepsilon_i\)。以上两步表明，任意导子都唯一通过发送 \(\varepsilon_i\mapsto m_i\) 的线性映射 \(E\to M\) 分解。所以 \((E,\delta)\) 满足微分模的泛性质，得到所需同构。最后，变量的导数可以同时任意指定，故导子空间对应直积；微分模的每个元素只有有限支集，故微分模本身对应直和。
\end{proof}

这个计算也说明有限型代数的微分模有限生成。若 \(B=A[b_1,\ldots,b_n]\)，每个元素均写作 \(f(b_1,\ldots,b_n)\)，相同的乘法法则给出
\[
	\mathrm{d}\bigl(f(b_1,\ldots,b_n)\bigr)
	=\sum_i\frac{\partial f}{\partial X_i}(b_1,\ldots,b_n)\,\mathrm{d}b_i.
\]
所以 \(\mathrm{d}b_1,\ldots,\mathrm{d}b_n\) 生成 \(\Omega_{B/A}\)。这些生成元是否独立，则取决于 \(b_i\) 满足的关系。

\begin{proposition}\label{g5:prop:localized-differentials}
	设 \(A\to B\) 为交换含幺环同态，\(S\subseteq B\) 为含 \(1\) 的乘法集，\(R=S^{-1}B\)。对任意 \(R\) 模 \(M\)，限制给出双射
	\[
		\operatorname{Der}_A(R,M)
		\xrightarrow{\ \sim\ }\operatorname{Der}_A(B,M).
	\]
	\(D:B\to M\) 的唯一延伸为
	\[
		\widetilde D(b/s)=s^{-1}D(b)-bs^{-2}D(s).
	\]
	这里 \(b,s\) 在系数位置表示其在 \(R\) 中的像。相应地，存在自然的 \(R\) 模同构
	\[
		S^{-1}\Omega_{B/A}
		\cong R\otimes_B\Omega_{B/A}
		\xrightarrow{\ \sim\ }\Omega_{R/A},
		\qquad
		\frac{\mathrm{d}_{B/A}b}{s}
		\longmapsto s^{-1}\mathrm{d}_{R/A}(b/1).
	\]
	在此识别下，
	\begin{equation}\label{g5:eq:localized-universal-derivative}
		\mathrm{d}(b/s)
		=\frac{\mathrm{d}b}{s}-\frac{b\,\mathrm{d}s}{s^2}.
	\end{equation}
\end{proposition}
\begin{proof}
	唯一性来自\eqref{g5:eq:derivative-inverse}：一个延伸在 \(s^{-1}\) 上必须取值 \(-s^{-2}D(s)\)，从而在 \(b/s\) 上也只能取陈述中的值。存在性的关键是分式表示不唯一，尤其 \(S\) 中可能含有零因子。

	若 \(b/s=b'/s'\)，则存在 \(u\in S\) 使
	\[
		u(s'b-sb')=0.
	\]
	令 \(c=s'b-sb'\)。在 \(R\) 中 \(c=0\)，所以它在 \(M\) 上的作用也为零。将原环中的等式 \(uc=0\) 用 \(D\) 求导，得到
	\(uD(c)+cD(u)=0\)。由于 \(u\) 在 \(M\) 上可逆，便有
	\[
		s'D(b)+bD(s')-sD(b')-b'D(s)=0.
	\]
	把这个等式除以 \(ss'\)，再用 \(b/s=b'/s'\)，恰得
	\[
		s^{-1}D(b)-bs^{-2}D(s)
		=(s')^{-1}D(b')-b'(s')^{-2}D(s').
	\]
	因此分式公式良定义。把和写成 \((bt+cs)/(st)\) 并使用原导子的乘法法则，便得加法性；把乘积写成 \(bc/(st)\)，则得到
	\[
		\widetilde D\!\left(\frac bs\frac ct\right)
		=\frac bs\,\widetilde D(c/t)+\frac ct\,\widetilde D(b/s).
	\]
	对来自 \(A\) 的元素，公式仍给零，所以 \(\widetilde D\) 是 \(A\) 导子，并在分母为 \(1\) 时恢复 \(D\)。若 \(0\in S\)，则 \(R=0\)、\(M=0\)，同一断言仍成立。

	为得到模同构，将 \(b\mapsto1\otimes\mathrm{d}_{B/A}b\) 看作取值于 \(R\otimes_B\Omega_{B/A}\) 的导子。刚证的延伸公式使它唯一延伸到 \(R\)。普遍性质给出
	\[
		\Omega_{R/A}\longrightarrow R\otimes_B\Omega_{B/A}.
	\]
	反方向由\eqref{g5:eq:differentials-natural-map}给出。两个复合分别在 \(1\otimes\mathrm{d}b\) 和 \(\mathrm{d}(b/s)\) 上为恒等，后一项使用已证分式公式；这些元素生成相应的模，所以两张映射互逆。最后用\cref{g5:prop:localization-tensor}识别张量积与模局部化，即得所述形式。
\end{proof}

例如，Laurent 多项式环 \(B=A[t,t^{-1}]\) 满足
\[
	\Omega_{B/A}=B\,\mathrm{d}t,\qquad
	\mathrm{d}(t^n)=nt^{n-1}\mathrm{d}t\quad(n\in\mathbb Z).
\]
若 \(k\) 为域，局部化进一步给出
\[
	\Omega_{k(x_1,\ldots,x_n)/k}
	=\bigoplus_{i=1}^n k(x_1,\ldots,x_n)\,\mathrm{d}x_i.
\]
此外，若先把基环中的一个乘法集 \(T\subseteq A\) 反演，则 \(A\) 导子延伸后自动杀死 \(T^{-1}A\)：它既杀死 \(a\)，也由逆元公式杀死 \(t^{-1}\)。所以
\[
	\Omega_{(T^{-1}B)/(T^{-1}A)}
	\cong T^{-1}\Omega_{B/A}.
\]
局部化中的求导公式也见\cite[第25节，习题25.3，第198页]{Matsumura1989}。

\ProseParagraphBreak

最后考虑两个 \(A\) 代数 \(B,C\) 的张量积 \(R=B\otimes_A C\)。其中
\(b\otimes c=(b\otimes1)(1\otimes c)\)，所以它的导数应当分别包含 \(B\) 方向与 \(C\) 方向的变化。两种变化可以独立指定，而公共基环 \(A\) 在两个方向都被当作常数。

\begin{proposition}\label{g5:prop:tensor-differentials}
	设 \(A\) 为交换含幺环，\(B,C\) 为交换含幺 \(A\) 代数，\(R=B\otimes_A C\)。存在自然的 \(R\) 模同构
	\[
		\Omega_{R/A}
		\cong
		\bigl(R\otimes_B\Omega_{B/A}\bigr)
		\oplus
		\bigl(R\otimes_C\Omega_{C/A}\bigr).
	\]
	在此同构下，
	\[
		\mathrm{d}(b\otimes c)
		\longmapsto
		\bigl((1\otimes c)\otimes\mathrm{d}b,\,
		(b\otimes1)\otimes\mathrm{d}c\bigr).
	\]
\end{proposition}
\begin{proof}
	对任意 \(R\) 模 \(M\)，一个 \(A\) 导子 \(D:R\to M\) 限制到两个自然映射 \(B\to R\) 与 \(C\to R\) 后，给出
	\[
		D_B(b)=D(b\otimes1),\qquad D_C(c)=D(1\otimes c).
	\]
	乘法法则强制
	\[
		D(b\otimes c)
		=(1\otimes c)D_B(b)+(b\otimes1)D_C(c).
	\]
	因此 \(D\) 由这两个限制唯一决定。

	反过来，任意给定 \(D_B\in\operatorname{Der}_A(B,M)\) 及
	\(D_C\in\operatorname{Der}_A(C,M)\)，用上式定义纯张量的值。右边对 \(b,c\) 都具有加法性，并尊重 \(A\) 的平衡关系：对 \(a\in A\)，导子在 \(a\) 上为零，所以把 \(a\) 乘到 \(b\) 上或乘到 \(c\) 上，都把右边乘以同一个 \(A\) 标量。因此它唯一给出加法映射 \(D:R\to M\)。在两个纯张量的乘积上，将 \(D_B(bb')\) 和 \(D_C(cc')\) 各自展开，所得四项恰为
	\[
		(b\otimes c)D(b'\otimes c')
		+(b'\otimes c')D(b\otimes c).
	\]
	加法展开使乘法法则对任意张量和成立。来自 \(A\) 的元素的导数为零，故 \(D\) 是所需的 \(A\) 导子。

	于是，取值于 \(M\) 的导子等价于独立选择上述两个导子。分别用两个微分模的泛性质和标量扩张的泛性质，这又等价于选择两个 \(R\) 线性映射
	\[
		R\otimes_B\Omega_{B/A}\longrightarrow M,\qquad
		R\otimes_C\Omega_{C/A}\longrightarrow M,
	\]
	也即从它们的直和到 \(M\) 的一个线性映射。陈述中的公式给出直和值导子，所以这个直和连同该导子满足 \(\Omega_{R/A}\) 的泛性质。唯一性便给出同构及其自然性。
\end{proof}

例如取 \(B=A[x,x^{-1}]\)、\(C=A[y]\)，张量积为 \(R=A[x,x^{-1},y]\)。两个方向给出
\[
	\Omega_{R/A}=R\,\mathrm{d}x\oplus R\,\mathrm{d}y,\qquad
	\mathrm{d}(x^{-1}y^2)
	=-x^{-2}y^2\,\mathrm{d}x+2x^{-1}y\,\mathrm{d}y.
\]
这个公式在任意特征中成立；若 \(2=0\)，第二项就是零。张量计算依靠的是导子的泛性质，不需要 \(B\) 或 \(C\) 在 \(A\) 上平坦。

\subsection{可分与不可分扩张的微分}
\label{g5:subsec:0501:04}

域扩张 \(K\subseteq L\) 的微分模是一个 \(L\) 向量空间。超越元可以自由求导，而代数元的导数受到最小多项式的约束。这个约束的强弱由最小多项式的形式导数决定；在正特征中，它可能根本不给出非零关系。

\begin{proposition}\label{g5:prop:simple-field-differentials}
	设 \(K\) 为域，\(\alpha\) 为 \(K\) 上的代数元，\(L=K(\alpha)\)，\(f(T)\in K[T]\) 为 \(\alpha\) 的首一最小多项式。则
	\[
		\Omega_{L/K}\cong
		L e\big/L\bigl(f'(\alpha)e\bigr),
		\qquad \mathrm{d}_{L/K}\alpha\longleftrightarrow[e].
	\]
	因而当 \(\alpha\) 在 \(K\) 上可分时，\(\Omega_{L/K}=0\)；当 \(\alpha\) 不可分时，\(\Omega_{L/K}=L\,\mathrm{d}\alpha\) 是一维 \(L\) 向量空间。
\end{proposition}
\begin{proof}
	这里 \(K[\alpha]\) 已是域，故 \(L=K[\alpha]\)，每个元素可写成 \(g(\alpha)\)。对任意 \(L\) 模 \(M\) 和 \(K\) 导子 \(D:L\to M\)，令 \(m=D(\alpha)\)。多项式求导法则给出
	\[
		D(g(\alpha))=g'(\alpha)m,\qquad
		0=D(f(\alpha))=f'(\alpha)m.
	\]
	所以 \(D\) 唯一由满足 \(f'(\alpha)m=0\) 的元素 \(m\) 决定。

	反过来，若 \(m\in M\) 满足这个等式，就尝试定义
	\(D(g(\alpha))=g'(\alpha)m\)。若 \(g(\alpha)=h(\alpha)\)，则 \(g-h=qf\) 对某个 \(q\in K[T]\) 成立，于是
	\[
		(g'-h')(\alpha)m
		=\bigl(q'(\alpha)f(\alpha)+q(\alpha)f'(\alpha)\bigr)m=0.
	\]
	故定义与多项式代表无关。多项式的加法、乘法求导公式又保证它为 \(K\) 导子。因此，给出一个导子恰等价于给出一个被 \(f'(\alpha)\) 杀死的 \(m\)，也就等价于给出一个线性映射
	\(Le/L(f'(\alpha)e)\to M\)。这个商模连同 \(g(\alpha)\mapsto g'(\alpha)[e]\) 满足普遍性质，得到公式。

	最后说明可分性为何恰好给出两种结果。不可约多项式 \(f\) 若 \(f'\ne0\)，则 \(\deg f'<\deg f\)，所以 \(\gcd(f,f')=1\)，从而没有重根，且 \(f'(\alpha)\ne0\)。这个非零元素在 \(L\) 中可逆，商模于是为零。若 \(f\) 不可分，则 \(f,f'\) 有公共因子；不可约性及次数比较迫使 \(f'=0\)，因而没有加在 \(e\) 上的非零关系。
\end{proof}

\begin{corollary}\label{g5:cor:separable-field-differentials}
	设 \(L/K\) 为可分代数域扩张，允许扩张次数无限。则
	\[
		\Omega_{L/K}=0,\qquad
		\operatorname{Der}_K(L,M)=0
	\]
	对每个 \(L\) 模 \(M\) 都成立。特别地，这适用于任意有限可分域扩张。
\end{corollary}
\begin{proof}
	任取 \(\alpha\in L\)，其在 \(K\) 上的最小多项式 \(f\) 可分。任意 \(K\) 导子 \(D:L\to M\) 都满足 \(f'(\alpha)D(\alpha)=0\)，而 \(f'(\alpha)\) 在 \(L\) 中可逆，故 \(D(\alpha)=0\)。每个 \(\alpha\) 都如此，所以 \(D=0\)。特别取普遍导子；其像生成 \(\Omega_{L/K}\)，故这个模也为零。
\end{proof}

例如 \(\mathbb Q(\sqrt2)/\mathbb Q\) 中，等式
\[
	0=\mathrm{d}\bigl((\sqrt2)^2-2\bigr)
	=2\sqrt2\,\mathrm{d}\sqrt2
\]
迫使 \(\mathrm{d}\sqrt2=0\)。正特征下也会出现相同的消失：有限域扩张 \(\mathbb F_p\subseteq\mathbb F_{p^m}\) 中，每个元素都满足 \(T^{p^m}-T=0\)，这个多项式的导数为 \(-1\)，故对该等式求导就给出 \(\mathrm{d}\alpha=0\)。可分代数扩张不产生新的相对微分方向。一般可分代数扩张的更强延拓性质见\cite[第25节，定理25.3，第195--196页]{Matsumura1989}；上面的论证直接得到此处所需的消失结论。

\begin{example}\label{g5:exm:inseparable-field-differentials}
	设 \(p\) 为素数，\(t\) 为超越元，取
	\[
		L=\mathbb F_p(t),\qquad K=\mathbb F_p(t^p)\subseteq L.
	\]
	则 \([L:K]=p\)。确实，\(\alpha=t\) 满足 \(T^p-t^p\in K[T]\)，故次数至多为 \(p\)。若
	\(\sum_{i=0}^{p-1}a_i(t^p)t^i=0\)，把 \(a_i\) 的分母清去后，各项中 \(t\) 的指数分别模 \(p\) 同余于不同的 \(i\)，不可能相消，所以所有 \(a_i=0\)，次数也至少为 \(p\)。

	最小多项式的导数为零，因此
	\[
		\Omega_{L/K}=L\,\mathrm{d}t,\qquad
		\operatorname{Der}_K(L,L)=L\,\partial_t.
	\]
	这里 \(\partial_t\) 先在 \(\mathbb F_p[t]\) 上定义，再按局部化公式延伸到 \(L\)；它杀死 \(t^p\)，也杀死 \(t^p\) 的任意有理函数，故确实为 \(K\) 导子。代数扩张的超越次数为零，微分模却有维数一，原因是特征 \(p\) 下的代数关系 \(t^p\in K\) 求导后成为 \(0=0\)。

	把同一个 \(L\) 改看成其自身的代数时，\(\Omega_{L/L}=0\)；把它看成 \(\mathbb F_p\) 的有理函数域时，又有 \(\Omega_{L/\mathbb F_p}=L\,\mathrm{d}t\)。这三种基环说明，微分模记录的是相对于指定常数的变化。
\end{example}

纯不可分扩张的次数也不能直接当作微分模的维数。例如
\[
	K=\mathbb F_p(t^{p^r})\subseteq L=\mathbb F_p(t),\qquad r\geq1.
\]
按指数模 \(p^r\) 分类的同一论证给出 \([L:K]=p^r\)，但最小多项式 \(T^{p^r}-t^{p^r}\) 的导数仍为零，所以 \(\Omega_{L/K}\) 始终只有一维。域扩张的次数计算元素的线性自由度，微分模则计算求导关系留下的一阶自由度；在不可分情形中，这两种计数差异尤其明显。微分基与正特征域扩张的系统联系见\cite[第26节，微分基及定理26.5，第201--202页]{Matsumura1989}。
% 纲要条目排为编号小节标题，并进入目录与书签；说明段仍保留于纲要注释中。

% !TeX root = ../../Book5.tex
% 纲要标题：### 5.2 余切正合列
% 纲要位置：几何学计划纲要.md，第 5196 行。
% 正文状态：初稿；尚未正式审核。

% 纲要内容开始
% 
% * 商环的余法序列
% * 复合映射的传递序列
% * 微分模与余法模的基变换
% * 余法映射的核
% 
% 纲要内容结束

\section{余切正合列}
\label{g5:sec:0502}

设一条曲线由方程 \(f(x,y)=0\) 给出。沿曲线求导应当满足
\[
	\frac{\partial f}{\partial x}\,\mathrm{d}x+
	\frac{\partial f}{\partial y}\,\mathrm{d}y=0.
\]
这里的系数属于曲线的坐标环，因而还要模去 \(f\)。要把这个计算用于多个方程、一般基环或复合映射，需要同时回答两个问题：这些微分关系是否已给出全部关系，关系之间又是否可能存在依赖。余切正合列回答前一个问题；它的左端核保留后一种依赖。以下两个正合列可对照 \cite[第25节，定理25.1、25.2，第193—195页]{Matsumura1989}。

\subsection{商环的余法序列}
\label{g5:subsec:0502:01}

设 \(A\to P\) 为环同态，\(I\) 为 \(P\) 的理想，\(B=P/I\)。微分从 \(P\) 降到 \(B\) 时，\(I\) 中的函数都变成零，其微分也必须消失。另一方面，对 \(f,g\in I\)，乘法公式给出
\(\mathrm{d}(fg)=f\,\mathrm{d}g+g\,\mathrm{d}f\)，张量到 \(B\) 后已自动为零。因此记录方程的一阶变化时，应先把 \(I^2\) 除掉。

\begin{definition}\label{g5:def:conormal-module}
	设 \(I\) 为环 \(P\) 的理想，\(B=P/I\)。\(B\) 模 \(I/I^2\) 称为这个商呈示的\term[yufa-mo]{余法模}[conormal module]。它的 \(B\) 模结构由
	\[
		(p+I)(f+I^2)=pf+I^2
	\]
	给出；改变 \(p\) 的代表只改变一个 \(I^2\) 中的元素。
\end{definition}

\begin{theorem}[余法正合列]\label{g5:thm:conormal-sequence}
	设 \(A\to P\) 为交换环同态，\(I\subseteq P\) 为理想，\(B=P/I\)。存在自然的 \(B\) 模正合列
	\[
		I/I^2\xrightarrow{\delta}
		\Omega_{P/A}\otimes_P B\xrightarrow{\beta}
		\Omega_{B/A}\longrightarrow0,
	\]
	其中
	\[
		\delta(f+I^2)=\mathrm{d}_{P/A}f\otimes1,\qquad
		\beta(\mathrm{d}_{P/A}p\otimes b)=b\,\mathrm{d}_{B/A}\bar p.
	\]
\end{theorem}
\begin{proof}
	先检验第一张映射。上面的乘积公式说明它在 \(I^2\) 上为零；对 \(p\in P,f\in I\)，有
	\[
		\mathrm{d}(pf)\otimes1
		=\bar p\,(\mathrm{d}f\otimes1),
	\]
	所以它确为 \(B\) 线性映射。第二张映射由导子 \(p\mapsto\mathrm{d}\bar p\) 的泛性质得到；\(\Omega_{B/A}\) 由全部 \(\mathrm{d}\bar p\) 生成，故 \(\beta\) 满射。显然 \(\beta\delta=0\)。

	为了证没有遗漏任何关系，令
	\[
		Q=(\Omega_{P/A}\otimes_P B)/\operatorname{im}\delta.
	\]
	把 \(\bar p\in B\) 送到 \(\mathrm{d}p\otimes1\) 在 \(Q\) 中的类。改变代表 \(p\) 所增加的是某个 \(\mathrm{d}f\otimes1\)、\(f\in I\)，已经被除去，所以这是良定义的 \(A\) 导子 \(B\to Q\)。泛性质给出 \(\Omega_{B/A}\to Q\)，而 \(\beta\) 给出反向映射。两复合分别固定生成元 \(\mathrm{d}\bar p\) 与 \(\mathrm{d}p\otimes1\) 的类，故互逆。这证明 \(\ker\beta=\operatorname{im}\delta\)。
\end{proof}

\begin{corollary}\label{g5:cor:differentials-presentation}
	设 \(A\) 为交换环，\(P=A[x_1,\ldots,x_n]\)，\(I=(f_1,\ldots,f_r)\)，\(B=P/I\)。则
	\[
		B^r\xrightarrow{J^{\mathsf T}}B^n
		\longrightarrow\Omega_{B/A}\longrightarrow0
	\]
	正合，其中第 \(i\) 个标准基向量送到
	\[
		\left(\overline{\frac{\partial f_i}{\partial x_1}},
		\ldots,\overline{\frac{\partial f_i}{\partial x_n}}\right).
	\]
	右端映射把 \(B^n\) 的第 \(j\) 个标准基向量送到 \(\mathrm{d}\bar x_j\)。
\end{corollary}
\begin{proof}
	由\cref{g5:prop:polynomial-differentials}，\(\Omega_{P/A}\) 自由于 \(\mathrm{d}x_1,\ldots,\mathrm{d}x_n\)。若 \(f=\sum_i a_if_i\in I\)，则
	\[
		\mathrm{d}f\otimes1=\sum_i\bar a_i(\mathrm{d}f_i\otimes1).
	\]
	因此余法映射的像由所列向量生成，定理即给出该呈示。
\end{proof}

这个公式适用于非约化环，也适用于正特征。它使用整个微分模，而非只在某一点把系数代成数。若换用另一组生成方程，得到的矩阵可能不同，但它们都给出由泛性质唯一确定的 \(\Omega_{B/A}\)。

\begin{example}\label{g5:exm:cusp-differential-torsion}
	设 \(k\) 为特征零的域，\(B=k[x,y]/(y^2-x^3)\)。用 \(x,y\) 同时表示商环中的坐标，则
	\[
		\Omega_{B/k}
		=\frac{B\,\mathrm{d}x\oplus B\,\mathrm{d}y}
		{B(-3x^2\,\mathrm{d}x+2y\,\mathrm{d}y)}.
	\]
	在正规化 \(B\subset k[t]\)、\(x=t^2,y=t^3\) 下，
	\(\mathrm{d}x\mapsto2t\,\mathrm{d}t\)，
	\(\mathrm{d}y\mapsto3t^2\,\mathrm{d}t\)。于是
	\(\tau=2x\,\mathrm{d}y-3y\,\mathrm{d}x\) 映到零，却在 \(\Omega_{B/k}\) 中非零。事实上，若
	\((-3y,2x)=c(-3x^2,2y)\)，则在分式域中
	\(c=y/x^2=t^{-1}\)，它不属于 \(B=k[t^2,t^3]\)。另一方面，利用微分关系可得
	\[
		y\tau=2xy\,\mathrm{d}y-3y^2\,\mathrm{d}x
		=3x^3\,\mathrm{d}x-3x^3\,\mathrm{d}x=0.
	\]
	所以整环的微分模也可能含有非零挠元。正规化映射在函数环上单射，并不迫使微分模上的映射单射。
\end{example}

\subsection{复合映射的传递序列}
\label{g5:subsec:0502:02}

现在设 \(A\to B\to C\)。把 \(C\) 中的函数对 \(A\) 求导时，\(B\) 中的元素也可以发生变化；改为对 \(B\) 求导，正好把这部分变化设为零。因而两种微分之间的商，由 \(B\) 相对于 \(A\) 的微分控制。

\begin{theorem}[传递正合列]\label{g5:thm:transitivity-sequence}
	设 \(A\to B\xrightarrow{u}C\) 为交换环同态。以下 \(C\) 模序列正合：
	\[
		C\otimes_B\Omega_{B/A}\xrightarrow{\alpha}
		\Omega_{C/A}\xrightarrow{\gamma}
		\Omega_{C/B}\longrightarrow0,
	\]
	其中
	\[
		\alpha(c\otimes\mathrm{d}_{B/A}b)=c\,\mathrm{d}_{C/A}u(b),
		\qquad
		\gamma(\mathrm{d}_{C/A}c)=\mathrm{d}_{C/B}c.
	\]
\end{theorem}
\begin{proof}
	\(b\mapsto\mathrm{d}_{C/A}u(b)\) 是值于 \(C\) 模的 \(A\) 导子，故先给出 \(B\) 线性映射，再延拓为 \(\alpha\)。\(\mathrm{d}_{C/B}\) 也为 \(A\) 导子，故给出 \(\gamma\)，并且 \(\gamma\) 满射、\(\gamma\alpha=0\)。

	令 \(Q=\Omega_{C/A}/\operatorname{im}\alpha\)。复合 \(C\to\Omega_{C/A}\to Q\) 在 \(u(B)\) 上为零，所以是 \(B\) 导子。它诱导 \(\Omega_{C/B}\to Q\)，与 \(\gamma\) 诱导的 \(Q\to\Omega_{C/B}\) 在全部微分生成元上互逆。于是中间正合。
\end{proof}

例如 \(C=B[t]\) 时，先把 \(\Omega_{B/A}\) 的系数扩到 \(C\)，再加上自由生成元 \(\mathrm{d}t\)，得到
\[
	\Omega_{B[t]/A}\cong
	(B[t]\otimes_B\Omega_{B/A})\oplus B[t]\,\mathrm{d}t.
\]
检验方法是给定一个 \(A\) 导子在 \(B\) 上的限制和在 \(t\) 上的任意值，用多项式乘法公式唯一延拓；因此上述直和满足微分模的泛性质。在这一例中传递序列左端单射，并且存在由多项式变量给出的分裂。

\begin{example}\label{g5:exm:inseparable-transitivity}
	设 \(k\) 为特征 \(p>0\) 的域，\(B=k[u]\)，\(C=k[t]\)，同态 \(B\to C\) 把 \(u\) 送到 \(t^p\)。这张环同态单射，而传递序列的第一张映射为
	\[
		C\,\mathrm{d}u\longrightarrow C\,\mathrm{d}t,
		\qquad \mathrm{d}u\longmapsto p t^{p-1}\mathrm{d}t=0.
	\]
	故 \(\Omega_{C/B}\cong C\,\mathrm{d}t\)。与之比较，若 \(u\mapsto t^n\) 其中 \(n\geq1\) 且 \(n\) 在 \(k\) 中非零，第一张映射为乘 \(nt^{n-1}\)，它在整环 \(C\) 上单射，而
	\(\Omega_{C/B}\cong C/(t^{n-1})\,\mathrm{d}t\)。当 \(n=1\) 时该商为零。
\end{example}

\subsection{微分模与余法模的基变换}
\label{g5:subsec:0502:03}

基变换使原来的系数环 \(A\) 变成 \(A'\)。在新环中，对 \(A'\) 求导时所有新系数都作为常数，因此只有原来 \(B\) 中相对于 \(A\) 的变化留下来。这个观察给出不需要平坦性假设的同构。

\begin{proposition}\label{g5:prop:differentials-base-change}
	设 \(A\to B\)、\(A\to A'\) 为交换环同态，\(B'=B\otimes_A A'\)。存在自然的 \(B'\) 模同构
	\[
		\Omega_{B/A}\otimes_B B'\xrightarrow{\ \sim\ }\Omega_{B'/A'},
		\qquad
		\mathrm{d}b\otimes(b_1\otimes a')
		\longmapsto (b_1\otimes a')\,\mathrm{d}(b\otimes1).
	\]
\end{proposition}
\begin{proof}
	记左端为 \(N\)。公式
	\[
		D(b\otimes a')=\mathrm{d}b\otimes(1\otimes a')
	\]
	对 \(A\) 平衡，因为 \(\mathrm{d}(ab)=a\,\mathrm{d}b\)。它由此给出加法映射 \(B'\to N\)。在纯张量乘积上直接应用 乘法公式可见 \(D\) 为 \(A'\) 导子；它在 \(1\otimes A'\) 上为零。因此得到 \(\Omega_{B'/A'}\to N\)。

	反向映射由 \(A\) 导子 \(b\mapsto\mathrm{d}(b\otimes1)\) 给出，再作标量扩张。两张映射的复合固定 \(\mathrm{d}b\otimes1\)，并且
	\[
		\mathrm{d}(b\otimes a')=(1\otimes a')\,\mathrm{d}(b\otimes1),
	\]
	说明也固定右端的全部生成元。因此互逆。
\end{proof}

在多项式商呈示中，这个同构就是把方程和偏导数矩阵的系数一同送到新基环。张量保持右正合，所以微分模的生成元与关系直接随基环改变。例如令
\[
	B=\mathbb Z[x]/(x^2-a),\qquad a\in\mathbb Z.
\]
则 \(\Omega_{B/\mathbb Z}=B/(2x)\,\mathrm{d}x\)。模二以后，
\[
	\Omega_{(B\otimes_{\mathbb Z}\mathbb F_2)/\mathbb F_2}
	=(B\otimes_{\mathbb Z}\mathbb F_2)\,\mathrm{d}x.
\]
导数关系消失，是系数 \(2\) 在新基环中变成零的结果。

\ProseParagraphBreak

余法模需要另作判断。若 \(B=P/I\)，记 \(P'=P\otimes_A A'\)，\(J=\operatorname{im}(I\otimes_A A'\to P')\)，则 \(B'=P'/J\)。有自然满射
\[
	(I/I^2)\otimes_B B'\longrightarrow J/J^2.
\]
其满射性来自 \(J\) 由 \(I\) 的像生成，乘积在两边都被除去；它未必单射。取 \(A=\mathbb Z\)、\(P=\mathbb Z[x]\)、\(I=(2x)\)、\(A'=\mathbb F_2\)。由于 \(2x\) 在 \(P\) 中不是零因子，映射 \(B\to I/I^2\)、\(b\mapsto2xb\) 为同构，故左端为 \(\mathbb F_2[x]\)；但 \(J=0\)，右端为零。这与微分模的基变换同构没有冲突：微分模由导子的泛性质控制，余法模还记录所选闭嵌入的方程。

\subsection{余法映射的核}
\label{g5:subsec:0502:04}

两个正合列的终点都为零，保证目标微分能由前面的生成元得到。起点没有写零，意味着方程的一阶变化之间可能还有核。这个核在一个方程的情形就已经出现。

\begin{example}\label{g5:exm:conormal-dual-number-kernel}
	设 \(k\) 为特征不为二的域，\(P=k[x]\)、\(I=(x^2)\)、\(B=k[x]/(x^2)\)。乘 \(x^2\) 给出 \(B\cong I/I^2\)，余法映射化为
	\[
		B\longrightarrow B\,\mathrm{d}x,\qquad
		b\longmapsto2bx\,\mathrm{d}x.
	\]
	它的核为 \((x)\)，其元素 \(x\) 对应 \(I/I^2\) 中非零的 \(x^3\) 类。因此
	\[
		\Omega_{B/k}\cong k\,\mathrm{d}x.
	\]
	若 \(\operatorname{char}k=2\)，这张余法映射处处为零，微分模则为自由 \(B\) 模 \(B\,\mathrm{d}x\)。两种情形中 \(x^2\) 都是多项式环中的非零因子，故方程在环境环中没有零因子，并不能保证其微分在商环上给出单射。
\end{example}

从导子的角度也能读出余法映射。对任意 \(B\) 模 \(M\)，由泛性质和商环下降条件得到正合列
\[
	0\longrightarrow\operatorname{Der}_A(B,M)
	\longrightarrow\operatorname{Der}_A(P,M)
	\longrightarrow\operatorname{Hom}_B(I/I^2,M),
\]
最后一张映射把 \(D\) 送到 \(f+I^2\mapsto D(f)\)。因为 \(I\) 在 \(M\) 上作用为零，乘法公式保证这确实是 \(B\) 线性映射；它的核恰为能降到 \(B\) 的导子。因此环境空间中的一阶运动保留商环方程，当且仅当它在余法模上的取值为零。

\ProseParagraphBreak

最后的映射也未必满射。例如上例取 \(M=k=B/(x)\)，每个导子都满足 \(D(x^2)=2xD(x)=0\)，而
\(\operatorname{Hom}_B(I/I^2,k)\cong k\) 非零。把方程的一阶变化任意指定为一个数，不一定能由环境坐标的一阶变化实现。下一节把 \(M\) 取为点的剩余域，将这一条件转成线性方程组，从而得到切空间的矩阵表示。

% !TeX root = ../../Book5.tex
% 纲要标题：### 5.3 切空间与 Jacobian 矩阵
% 纲要位置：几何学计划纲要.md，第 5203 行。
% 正文状态：初稿；尚未正式审核。

% 纲要内容开始
% 
% * 双数点及切向量
% * 极大理想模平方
% * Jacobian 矩阵的秩计算
% * 正则局部环的维数定义
% * 非完美域上的正则性与几何正则性
% 
% 纲要内容结束

\section{切空间与 Jacobian 矩阵}
\label{g5:sec:0503}

在方程 \(f(x_1,\ldots,x_n)=0\) 上选定一点以后，把坐标改为 \(x_i+\varepsilon v_i\)，并只保留 \(\varepsilon\) 的一次项，就得到切向量满足的线性条件。双数环把“只保留一次项”变成一个准确的环运算；微分模则把这个运算写成与坐标无关的形式。由此可以比较两种局部信息：素链测量的 Krull 维数，以及一阶函数所需的独立坐标数。二者相等时，局部环称为正则。基域与点的剩余域不同时，还须区分局部环自身的余切空间和相对于基域的微分；这一差别在非完美域上尤为明显。

\subsection{双数点及切向量}
\label{g5:subsec:0503:01}

设 \(K\) 为域。环
\[
	K[\varepsilon]/(\varepsilon^2)
	=\{a+\varepsilon b:a,b\in K\}
\]
称为 \(K\) 上的\term[shuangshuhuan]{双数环}[ring of dual numbers]，其中 \(\varepsilon\) 非零而 \(\varepsilon^2=0\)。投影 \(a+\varepsilon b\mapsto a\) 留下原来的点，系数 \(b\) 记录一次变化。这里没有极限过程，因而同样适用于正特征。

\begin{definition}\label{g5:def:dual-number-tangent}
	设 \(k\) 为域，\(A\) 为 \(k\) 代数，\(K/k\) 为域扩张，\(x:A\to K\) 为固定的 \(k\) 代数同态。若 \(k\) 代数同态
	\[
		\widetilde x:A\longrightarrow K[\varepsilon]/(\varepsilon^2)
	\]
	模去 \(\varepsilon\) 后得到 \(x\)，则称 \(\widetilde x\) 为 \(x\) 处相对于 \(k\) 的双数切向量。全部这样的提升构成的集合记为 \(\mathrm{T}_x(A/k)\)。
\end{definition}

每个提升都唯一写成 \(\widetilde x(a)=x(a)+\varepsilon D(a)\)。同态保持乘法，正好要求 \(D\) 满足导子的乘法法则；这使上述集合自然成为向量空间。

\begin{proposition}\label{g5:prop:tangent-derivations}
	设 \(k\) 为域，\(A\) 为 \(k\) 代数，\(K/k\) 为域扩张，\(x:A\to K\) 为 \(k\) 代数同态，并用 \(x\) 将 \(K\) 看作 \(A\) 模。取 \(\varepsilon\) 的系数给出自然对应
	\[
		\mathrm{T}_x(A/k)\cong\operatorname{Der}_k(A,K)
		\cong\operatorname{Hom}_K
		\bigl(\Omega_{A/k}\otimes_{A,x}K,K\bigr).
	\]
	由此赋予 \(\mathrm{T}_x(A/k)\) 一个 \(K\) 向量空间结构，其零向量是复合 \(A\xrightarrow{x}K\hookrightarrow K[\varepsilon]/(\varepsilon^2)\)。
\end{proposition}
\begin{proof}
	保持加法和 \(k\) 中的常数，分别等价于 \(D\) 可加且在 \(k\) 上为零。比较
	\[
		\widetilde x(a)\widetilde x(b)
		=x(a)x(b)+\varepsilon\bigl(x(a)D(b)+x(b)D(a)\bigr)
	\]
	与 \(\widetilde x(ab)\)，便知保持乘法等价于
	\[
		D(ab)=x(a)D(b)+x(b)D(a).
	\]
	反过来，每个这样的导子给出环同态；导子满足 \(D(1)=0\)，所以它也保持单位。第一项对应因而成立。

	由微分模的泛性质，\(\operatorname{Der}_k(A,K)\cong\operatorname{Hom}_A(\Omega_{A/k},K)\)。一个 \(A\) 线性映射 \(h:\Omega_{A/k}\to K\) 唯一延伸为 \(\omega\otimes c\mapsto c h(\omega)\) 的 \(K\) 线性映射，得到第二项对应。最后通过导子的加法和数乘定义切向量的线性运算；这里并不是把两个环同态直接相加。
\end{proof}

若 \(\varphi:A\to B\) 为 \(k\) 代数同态，\(y:B\to K\)，\(x=y\varphi\)，则复合给出
\[
	\mathrm{T}_y(B/k)\longrightarrow \mathrm{T}_x(A/k),\qquad
	D\longmapsto D\circ\varphi.
\]
这就是相应几何映射的切映射。它与复合相容，也是自然微分映射在点上取值后所得线性映射的对偶。因而“函数反向拉回，切向量正向映出”的方向已经包含在导子的复合公式中。

\subsection{极大理想模平方}
\label{g5:subsec:0503:02}

设 \((R,\mathfrak m,\kappa)\) 为局部环。\(\mathfrak m\) 中的函数在该点为零，\(\mathfrak m^2\) 则由两项这样的函数的乘积组成。因此 \(\mathfrak m/\mathfrak m^2\) 保留消失函数的一阶部分，称为该局部环的\term[yuqiekongjian]{余切空间}[cotangent space]。它自然是 \(\kappa=R/\mathfrak m\) 向量空间，因为 \(\mathfrak m\) 乘上任意元素都变为零。若点为 \(k\) 有理点，这个局部定义正好给出上一小节相对切空间的对偶。

\begin{theorem}\label{g5:thm:rational-cotangent-space}
	设 \(k\) 为域，\(A\) 为 \(k\) 代数，\(x:A\to k\) 为 \(k\) 代数同态。令 \(\mathfrak p=\ker x\)、\(R=A_{\mathfrak p}\)、\(\mathfrak m=\mathfrak pR\)，则剩余域由 \(x\) 识别为 \(k\)，并有自然同构
	\[
		\mathfrak p/\mathfrak p^2
		\cong\mathfrak m/\mathfrak m^2
		\xrightarrow{\ \sim\ }\Omega_{A/k}\otimes_A k.
	\]
	最后一张映射把消失函数的类送到它的微分在该点的值。特别地，
	\[
		\mathrm{T}_x(A/k)\cong
		\operatorname{Hom}_k(\mathfrak m/\mathfrak m^2,k).
	\]
\end{theorem}
\begin{proof}
	\(\mathfrak p/\mathfrak p^2\) 被 \(\mathfrak p\) 零化，故为 \(k\) 向量空间；每个 \(s\notin\mathfrak p\) 的作用是乘非零标量 \(x(s)\)。局部化因而不改变这个模，结合局部化的正合性得到第一项同构。由\cref{g5:prop:localized-differentials}，又有
	\[
		\Omega_{A/k}\otimes_A k
		\cong\Omega_{R/k}\otimes_R k.
	\]
	我们在局部环 \(R\) 上构造所需同构。

	映射 \(u\mapsto\mathrm{d}u\otimes1\)、\(u\in\mathfrak m\)，在 \(\mathfrak m^2\) 上为零，因为
	\(\mathrm{d}(uv)=u\,\mathrm{d}v+v\,\mathrm{d}u\)。它因此给出 \(k\) 线性映射
	\[
		\alpha:\mathfrak m/\mathfrak m^2
		\longrightarrow\Omega_{R/k}\otimes_R k.
	\]
	反向映射利用有理点条件：每个 \(r\in R\) 唯一写成 \(c+u\)，其中 \(c\in k\)、\(u\in\mathfrak m\)。令
	\[
		\Delta(r)=[r-\bar r]=[u]\in\mathfrak m/\mathfrak m^2,
	\]
	这里把剩余类 \(\bar r\in k\) 当作 \(R\) 中的常数。若 \(r=c+u\)、\(s=d+v\)，则 \(uv\) 的类为零，故
	\[
		\Delta(rs)=c[v]+d[u]
		=\bar r\,\Delta(s)+\bar s\,\Delta(r).
	\]
	所以 \(\Delta\) 是值于 \(\mathfrak m/\mathfrak m^2\) 的 \(k\) 导子，诱导反向映射 \(\beta\)。对 \(u\in\mathfrak m\)，有 \(\beta\alpha([u])=[u]\)；对任意 \(r\in R\)，有
	\[
		\alpha\beta(\mathrm{d}r\otimes1)
		=\mathrm{d}(r-\bar r)\otimes1
		=\mathrm{d}r\otimes1.
	\]
	这些元素生成相应空间，故两映射互逆。最后应用\cref{g5:prop:tangent-derivations}即得切空间的对偶描述。
\end{proof}

上述同构不需要选择局部坐标，但用到了剩余域就是基域。一般点没有从剩余域到局部环的指定截面，\(r-\bar r\) 便不能按同样方式理解。此时应保留完整的余切正合尾。

\begin{proposition}\label{g5:prop:residue-cotangent-tail}
	设 \(k\) 为域，\(A\) 为 \(k\) 代数，\(\mathfrak p\in\operatorname{Spec}A\)。记 \(R=A_{\mathfrak p}\)、\(\mathfrak m=\mathfrak pR\)、\(\kappa=R/\mathfrak m\)，并令 \(x:A\to\kappa\) 为自然同态。则有 \(\kappa\) 向量空间正合列
	\[
		\mathfrak m/\mathfrak m^2
		\xrightarrow{\alpha}\Omega_{A/k}\otimes_A\kappa
		\longrightarrow\Omega_{\kappa/k}\longrightarrow0,
	\]
	其中 \(\alpha([u])=\mathrm{d}u\otimes1\)，右侧使用微分与局部化的自然识别。对偶后有正合列
	\[
		0\longrightarrow\operatorname{Der}_k(\kappa,\kappa)
		\longrightarrow \mathrm{T}_x(A/k)
		\longrightarrow
		\operatorname{Hom}_\kappa(\mathfrak m/\mathfrak m^2,\kappa).
	\]
	第一行的 \(\alpha\) 未必单射，第二行最右侧的映射未必满射。
\end{proposition}
\begin{proof}
	对商同态 \(R\to\kappa=R/\mathfrak m\) 应用\cref{g5:thm:conormal-sequence}，再以\cref{g5:prop:localized-differentials}识别中间微分模，便得到第一行。向量空间的线性泛函可以从子空间延伸到全空间，因此取 \(\kappa\) 线性对偶保持正合。结合导子的泛性质得到第二行；它的第一张映射是与 \(A\to\kappa\) 复合，最后一张映射则是把导子唯一延伸到 \(R\)，再限制于 \(\mathfrak m\) 并模去 \(\mathfrak m^2\)。
\end{proof}

例如取 \(A=k[t]\) 及素理想 \((0)\)。此时 \(R=\kappa=k(t)\)、\(\mathfrak m=0\)，局部余切空间为零；但
\[
	\Omega_{A/k}\otimes_A\kappa
	=\kappa\,\mathrm{d}t=\Omega_{\kappa/k}.
\]
相对于 \(k\) 的切向量仍有一维，它来自剩余域中 \(t\) 的变化。这已经说明，离开有理点后不能把两种余切空间直接等同。最后一小节还会给出闭点处 \(\alpha\) 为零而源非零的例子。

\subsection{Jacobian 矩阵的秩计算}
\label{g5:subsec:0503:03}

现在用方程计算切空间。设 \(A=k[X_1,\ldots,X_n]/(f_1,\ldots,f_r)\)，选定 \(K\) 值点 \(x:A\to K\)，记 \(a_j=x(\bar X_j)\)。由偏导数组成的矩阵
\[
	J_x=\left(\frac{\partial f_i}{\partial X_j}(a_1,\ldots,a_n)\right)
	_{\substack{1\leq i\leq r\\1\leq j\leq n}}
\]
称为这组方程在 \(x\) 处的\term[Jacobianjuzhen]{Jacobian 矩阵}[Jacobian matrix]。
\footnote{Carl Gustav Jacob Jacobi（1804--1851），德国数学家，研究椭圆函数、偏微分方程及函数行列式。}它把坐标的一次变化送到方程的一次变化，所以保留全部方程的切向量应当位于它的核。

\begin{theorem}\label{g5:thm:jacobian-tangent-space}
	设 \(k\) 为域，\(A=k[X_1,\ldots,X_n]/(f_1,\ldots,f_r)\)，\(K/k\) 为域扩张，\(x:A\to K\) 为 \(k\) 代数同态。按上述记号，取各坐标的导数值给出同构
	\[
		\mathrm{T}_x(A/k)\xrightarrow{\ \sim\ }
		\ker(J_x:K^n\to K^r),\qquad
		D\longmapsto\bigl(D(\bar X_1),\ldots,D(\bar X_n)\bigr).
	\]
	特别地，
	\[
		\dim_K \mathrm{T}_x(A/k)=n-\operatorname{rank}_K J_x.
	\]
\end{theorem}
\begin{proof}
	给定 \(v=(v_1,\ldots,v_n)\in K^n\)，存在唯一 \(k\) 代数同态
	\[
		k[X_1,\ldots,X_n]\longrightarrow K[\varepsilon]/(\varepsilon^2),
		\qquad X_j\longmapsto a_j+\varepsilon v_j.
	\]
	在单项式上展开乘积，所有含两个或更多 \(\varepsilon\) 的项均为零，得到对任意多项式 \(f\) 都成立的等式
	\[
		f(a+\varepsilon v)
		=f(a)+\varepsilon\sum_{j=1}^n
		\frac{\partial f}{\partial X_j}(a)v_j.
	\]
	这个展开只使用整数系数的二项式公式，在任意特征下都成立。由于 \(f_i(a)=0\)，该同态能通过商环 \(A\)，当且仅当每个 \(f_i\) 的一次项为零，即 \(J_xv=0\)。由切向量与导子的对应得到所需同构，维数公式再由秩零度定理得到。
\end{proof}

同一计算也可以从微分呈示读出。把\cref{g5:cor:differentials-presentation}张量到 \(K\)，得到
\[
	K^r\xrightarrow{J_x^{\mathsf T}}K^n
	\longrightarrow\Omega_{A/k}\otimes_AK\longrightarrow0.
\]
取对偶后，核中的矩阵便是 \(J_x\)。这说明 Jacobian 核与微分模对偶描述的是同一个空间。若 \(g=\sum_i h_if_i\)，则
\[
	\mathrm{d}g(a)=\sum_i h_i(a)\,\mathrm{d}f_i(a),
\]
因为含 \(f_i(a)\) 的项均为零。因此更换同一理想的生成元不会改变这些微分的行空间；更换整个多项式呈示时，内在定义 \(\operatorname{Der}_k(A,K)\) 仍给出各个核之间的自然同构。

\begin{example}\label{g5:exm:parabola-tangent}
	设 \(k\) 为任意域，\(A=k[x,y]/(y-x^2)\)。在有理点 \((a,a^2)\)，Jacobian 行为 \((-2a,1)\)，所以切向量恰为
	\[
		(v,2av),\qquad v\in k.
	\]
	参数 \(t\mapsto(t,t^2)\) 的切映射也把 \(v\) 送到这个向量。在特征二时，第二个分量恒为零，但切空间仍为一维；秩始终为一，因为对 \(y\) 的偏导等于 \(1\)。
\end{example}

\begin{example}\label{g5:exm:double-point-tangent}
	设 \(A=k[z]/(z^2)\)，在唯一的有理点 \(z=0\) 计算。方程的偏导 \(2z\) 在点上为零，所以 \(\mathrm{T}_0(A/k)\cong k\)，与特征无关。每个 \(v\in k\) 都给出提升 \(z\mapsto\varepsilon v\)。约化商 \(A/(z)=k\) 的相对切空间却为零。因而切空间能检测一些在底层点集中看不见的幂零结构。
\end{example}

\subsection{正则局部环的维数定义}
\label{g5:subsec:0503:04}

对非零 Noether 局部环 \((R,\mathfrak m,\kappa)\)，第三章已经定义嵌入维数
\[
	\operatorname{edim}R=\dim_\kappa\mathfrak m/\mathfrak m^2
\]
并证明 \(\dim R\leq\operatorname{edim}R<\infty\)，见\cref{g5:prop:dimension-embedding-bound}。Nakayama 引理把右端解释为生成 \(\mathfrak m\) 所需的最少元素数；高度定理说明它不能小于 Krull 维数。正则性要求这个下界达到。

\begin{definition}\label{g5:def:regular-local-ring}
	设 \((R,\mathfrak m,\kappa)\) 为非零 Noether 局部环。若
	\[
		\operatorname{edim}R=\dim R,
	\]
	则称 \(R\) 为\term[zhengzejubuhuan]{正则局部环}[regular local ring]。对 Noether 环 \(A\)，若每个素理想 \(\mathfrak p\) 处的局部环 \(A_{\mathfrak p}\) 都正则，则称 \(A\) 为\term[zhengzehuan]{正则环}[regular ring]。
\end{definition}

等价地，\(d\) 维 Noether 局部环正则，当且仅当它的极大理想能由 \(d\) 个元素生成。事实上，若余切空间有 \(d\) 维，提升一组基后由\cref{g5:cor:nakayama-generators}生成极大理想；反过来，\(d\) 个生成元使嵌入维数至多为 \(d\)，再与维数下界合用便得等号。这个定义和生成元解释可对照\cite[第14节开头，第105页]{Matsumura1989}。

\begin{example}\label{g5:exm:basic-regular-local-rings}
	每个域都是零维正则局部环。反过来，零维 Noether 局部环若正则，就有 \(\mathfrak m/\mathfrak m^2=0\)；对有限模 \(\mathfrak m\) 用 Nakayama 引理得到 \(\mathfrak m=0\)，故它是域。因此双数环 \(k[z]/(z^2)\) 虽为零维局部环，却不正则。

	离散赋值环的维数为一，极大理想由一个非零元素生成，所以也是正则局部环。另取 \(k\) 有理点 \(a=(a_1,\ldots,a_n)\)，多项式局部环
	\[
		S=k[X_1,\ldots,X_n]_{(X_1-a_1,\ldots,X_n-a_n)}
	\]
	的维数为 \(n\)：依次添加 \(X_j-a_j\) 给出长度为 \(n\) 的严格素链，而\cref{g5:cor:polynomial-dimension}给出上界。它的极大理想恰可由这 \(n\) 个元素生成，故 \(S\) 正则。
\end{example}

\begin{corollary}[有理点处的 Jacobian 正则性判别]\label{g5:cor:rational-jacobian-regularity}
	设 \(k\) 为域，\(A=k[X_1,\ldots,X_n]/(f_1,\ldots,f_r)\)，\(a\in k^n\) 为方程的公共零点。令 \(\mathfrak m\) 为该点的极大理想，\(d=\dim A_{\mathfrak m}\)。则
	\[
		A_{\mathfrak m}\text{ 正则}
		\quad\Longleftrightarrow\quad
		\operatorname{rank}_kJ_a=n-d.
	\]
\end{corollary}
\begin{proof}
	\(A\) 是有限型 \(k\) 代数，故 \(A_{\mathfrak m}\) 为 Noether 局部环。由\cref{g5:thm:rational-cotangent-space,g5:thm:jacobian-tangent-space}，
	\[
		\operatorname{edim}A_{\mathfrak m}
		=\dim_k \mathrm{T}_a(A/k)=n-\operatorname{rank}_kJ_a.
	\]
	把它代入正则局部环的定义即得结论。
\end{proof}

这里的 \(d\) 是该点的局部维数。这个判别不要求 \(k\) 代数闭，也不要求 \(A\) 约化或等维。经典仿射代数集的对应叙述及局部余切空间计算，见\cite[第I章第5节，定理5.1，第31—32页]{Hartshorne1977}。

\begin{example}\label{g5:exm:hypersurface-rational-regularity}
	设 \(f\in k[X_1,\ldots,X_n]\) 非恒定，\(a\in k^n\) 满足 \(f(a)=0\)。则 \(A=k[X_1,\ldots,X_n]/(f)\) 在该点的局部维数为 \(n-1\)。为说明这一步，令 \(S\) 为上例中的多项式局部环。因为 \(f\in\mathfrak m_S\)，\cref{g5:lem:principal-dimension-drop}给出
	\[
		\dim S/(f)\geq\dim S-1=n-1.
	\]
	另一方面，\cref{g5:prop:polynomial-hypersurface-dimension}给出 \(\dim A=n-1\)，局部化后维数不增，故得到等号。于是该点正则，当且仅当
	\[
		\left(\frac{\partial f}{\partial X_1}(a),\ldots,
		\frac{\partial f}{\partial X_n}(a)\right)\ne0.
	\]
	例如尖点 \(y^2-x^3=0\) 在原点的两项偏导都为零，因此局部维数为一而嵌入维数为二。结点 \(y^2-x^2(x+1)=0\)、\(\operatorname{char}k\ne2\)，在原点也有同样的两种维数。二者都不正则；仅靠切空间维数还不能区别这两种曲线奇点，下一节的切锥将保留更细的信息。
\end{example}

\begin{example}\label{g5:exm:jacobian-local-dimension}
	考虑\cref{g5:exm:mixed-dimensional-components}中的
	\[
		A=k[x,y,z]/(xy,xz).
	\]
	它有一个二维平面分支和一个一维直线分支，故 \(\dim A=2\)。在有理点 \(a=(1,0,0)\)，Jacobian 矩阵为
	\[
		J_a=\begin{pmatrix}0&1&0\\0&0&1\end{pmatrix},
	\]
	切空间维数为一。在该点附近反演 \(x\)，两个方程强迫 \(y=z=0\)，所以局部环为 \(k[x]_{(x-1)}\)，维数为一且正则。若把判别中的局部维数误换为全环维数二，就会得到错误判断。
\end{example}

\subsection{非完美域上的正则性与几何正则性}
\label{g5:subsec:0503:05}

正则性只检查给定环的局部环。若环带有 \(k\) 代数结构，还可以把系数扩张到更大的域，再检查正则性；这提出了额外要求。我们先记录这个要求，再用一个可直接计算的例子说明其作用。

\begin{definition}\label{g5:def:geometrically-regular-algebra}
	设 \(k\) 为域，\(A\) 为 Noether \(k\) 代数。若对每个有限域扩张 \(K/k\)，环 \(A\otimes_kK\) 都正则，则称 \(A\) 在 \(k\) 上\term[jihezhengze]{几何正则}[geometrically regular]。
\end{definition}

有限域扩张下 \(A\otimes_kK\) 是有限 \(A\) 模，因此仍为 Noether 环。定义中的检验确有意义；取 \(K=k\) 又可知几何正则必蕴含正则。这一有限域扩张版本的定义见\cite[第28节，第219页]{Matsumura1989}。

\begin{example}\label{g5:exm:inseparable-regular-point}
	设 \(p\) 为素数，\(k=\mathbb F_p(s)\)，其中 \(s\) 是超越元，并令
	\[
		L=k[\alpha]=k[T]/(T^p-s),\qquad\alpha^p=s.
	\]
	这是一个次数为 \(p\) 的域扩张。可以直接把它实现为 \(\mathbb F_p(u^p)\subset\mathbb F_p(u)\)，令 \(s=u^p\)、\(\alpha=u\)：\(1,u,\ldots,u^{p-1}\) 在 \(\mathbb F_p(u^p)\) 上线性无关，因为清去分母后不同项的指数模 \(p\) 不同；它们也张成整个域，因为对非零多项式 \(q(u)\)，有
	\[
		\frac1{q(u)}=\frac{q(u)^{p-1}}{q(u)^p},
		\qquad q(u)^p\in\mathbb F_p(u^p).
	\]
	先将任意多项式按指数模 \(p\) 分组，便得到张成性。由此既证明该商是域，也说明 \(s\) 在 \(k\) 中没有 \(p\) 次根，所以 \(k\) 不是完美域。

	作为域，\(L\) 是零维正则局部环，其极大理想为零。但由微分呈示，
	\[
		\Omega_{L/k}=L\,\mathrm{d}\alpha,
	\]
	因为关系的微分 \(\mathrm{d}(T^p-s)\) 为零。因此在恒等点 \(L\to L\) 处，相对于 \(k\) 的双数切空间有一维，而局部环的余切空间为零。这里\cref{g5:prop:residue-cotangent-tail}的全部非零信息都来自剩余域的微分。

	现在作有限基扩张 \(k\subset L\)。以第二个 \(L\) 为系数域，有
	\[
		L\otimes_kL
		\cong L[T]/(T^p-s)
		\cong L[\eta]/(\eta^p),\qquad\eta=T-\alpha.
	\]
	右端只有一个素理想 \((\eta)\)，因而维数为零；其极大理想模平方却有一维，由 \(\eta\) 的类生成。因此它不正则。原来的域 \(L\) 正则，却不在 \(k\) 上几何正则。
\end{example}

这个例子还可以放回一条正则曲线上的非有理闭点。取 \(A=k[z]\)、\(\mathfrak p=(z^p-s)\)，则剩余域为 \(L\)，\(R=A_{\mathfrak p}\) 是主理想整环在非零素理想处的局部化，极大理想由 \(f=z^p-s\) 生成。它的素理想只有 \((0)\) 和 \(fR\)，故 \(\dim R=1\)，且 \(R\) 正则。类 \([f]\) 生成一维 \(L\) 空间 \(\mathfrak m/\mathfrak m^2\)：它非零，因为若 \(f\in f^2R\)，在整环中约去 \(f\) 会使 \(f\) 成为单位。但余切正合尾为
\[
	L\cdot[f]\xrightarrow{\ 0\ }L\,\mathrm{d}z
	\xrightarrow{\ \sim\ }L\,\mathrm{d}\alpha\longrightarrow0.
\]
第一张映射为零正是因为 \(\mathrm{d}f=0\)。所以即使局部环本身正则，也不能省掉剩余域微分，并把余法映射默认成单射。

\ProseParagraphBreak

还可以直接检验这个域扩张的一次提升障碍。仍令 \(f=T^p-s\)，取
\[
	C=k[T]/(f^2),\qquad N=(f)/(f^2)\subset C.
\]
则 \(N^2=0\) 且 \(C/N=L\)。恒等同态 \(L\to C/N\) 不能提升为 \(k\) 代数同态 \(L\to C\)。事实上，任何提升都必须把 \(\alpha\) 送到某个 \(T+fh\)，其中 \(h\in C\)。在特征 \(p\) 下，
\[
	(T+fh)^p-s=f+f^ph^p=f\quad\text{在 }C\text{ 中},
\]
因为 \(p\geq2\) 且 \(f^2=0\)。但 \(f\) 在 \(C\) 中非零，这与 \(\alpha^p=s\) 矛盾。因而 \(L\) 本身的正则性并不保证这种平方零理想上的提升能够完成。

\ProseParagraphBreak

后文将在\secref{g5:sec:1501}用提升问题定义形式光滑等相对性质，并在\secref{g5:sec:1502}建立光滑态射的结构和相对 Jacobian 判据。本节已经证明的有理点判别针对局部环的正则性；对一般剩余域以及基域扩张，还必须保留相对微分与几何正则性的条件。上面的纯不可分计算给出了这一区别的具体原因。
% 纲要条目排为编号小节标题，并进入目录与书签；说明段仍保留于纲要注释中。

% !TeX root = ../../Book5.tex
% 纲要标题：### 5.4 切锥与高阶局部信息
% 纲要位置：几何学计划纲要.md，第 5211 行。
% 正文状态：初稿；尚未正式审核。

% 纲要内容开始
% 
% * 伴随分次环与最低次初始理想
% * 切锥与切空间
% * 尖点、结点和加厚点
% * 形式同构与有限阶截断的差异
% 
% ---
% 
% 纲要内容结束

\section{切锥与高阶局部信息}
\label{g5:sec:0504}

尖点与结点在原点的一阶方程都为零，所以它们的切空间都是整个平面；但尖点只有一个切线方向，结点在特征不为二时有两个。这些差别藏在一阶之后。把局部函数按消失阶逐层排列，可以得到切锥；再保留所有阶之间相容的关系，就进入形式方程。本节分别建立这两种构造，并说明它们各自保存什么信息。

\subsection{伴随分次环与最低次初始理想}
\label{g5:subsec:0504:01}

设 \((R,\mathfrak m,\kappa)\) 为非零 Noether 局部环。商 \(\mathfrak m/\mathfrak m^2\) 保留一次变化，下一层 \(\mathfrak m^2/\mathfrak m^3\) 保留无法由更高阶项改变的二次信息。理想幂满足 \(\mathfrak m^i\mathfrak m^j\subseteq\mathfrak m^{i+j}\)，所以这些相邻商还可以彼此相乘。

\begin{definition}
\label{g5:def:associated-graded}
	设 \((R,\mathfrak m,\kappa)\) 为非零 Noether 局部环。约定 \(\mathfrak m^0=R\)，定义其关于极大理想的\term[bansuifencihuan]{伴随分次环}[associated graded ring]为
	\[
		\operatorname{gr}_{\mathfrak m}R
		=\bigoplus_{d\geq0}\mathfrak m^d/\mathfrak m^{d+1},
	\]
	其中乘法由
	\[
		(a+\mathfrak m^{i+1})(b+\mathfrak m^{j+1})
		=ab+\mathfrak m^{i+j+1}
	\]
	给出。第 \(d\) 项称为次数 \(d\) 的齐次部分。
\end{definition}
\symidx{\(\operatorname{gr}_{\mathfrak m}R\)：局部环的伴随分次环}

若把 \(a\) 改为 \(a+a'\)、\(a'\in\mathfrak m^{i+1}\)，乘积的变化属于 \(\mathfrak m^{i+j+1}\)；改变 \(b\) 时同理，所以乘法良定义。零次部分是 \(\kappa=R/\mathfrak m\)，它使每个齐次部分成为 \(\kappa\) 向量空间。这里不需要在 \(R\) 中选出一个与 \(\kappa\) 同构的子域：标量在 \(R\) 中的两个提升相差 \(\mathfrak m\)，作用在第 \(d\) 层上相差的项自然落入 \(\mathfrak m^{d+1}\)。这一构造可对照\cite[第13节，第93页]{Matsumura1989}。

\ProseParagraphBreak

要从局部方程算出这个环，应当取多项式的最低次部分。设 \(P=k[x_1,\ldots,x_n]\)，\(\mathfrak n=(x_1,\ldots,x_n)\)。将非零多项式写成
\[
	f=f_d+f_{d+1}+\cdots,\qquad f_d\ne0,
\]
其中 \(f_j\) 为 \(j\) 次齐次多项式。记 \(\operatorname{ord}_0(f)=d\)，并称 \(f^\ast=f_d\) 为它在原点的\term[chushixingshi]{初始形式}[initial form]。对理想 \(I\subseteq\mathfrak n\)，定义
\[
	\operatorname{in}_0(I)
	=\bigl(f^\ast:0\ne f\in I\bigr)\subseteq k[X_1,\ldots,X_n],
\]
其中把初始形式中的 \(x_i\) 写成 \(X_i\)，强调它是分次环中的变量。这称为最低次初始理想；其生成元取自整个理想 \(I\)。第一章单项式序的初始理想取排序后的首项，这里取的是按消失阶排列的最低次齐次部分，两种构造的排序方式不同。
\symidx{\(\operatorname{in}_0(I)\)：原点处的最低次初始理想}

\begin{proposition}
\label{g5:prop:initial-ideal-tangent-cone}
	设 \(k\) 为域，\(P=k[x_1,\ldots,x_n]\)，\(\mathfrak n=(x_1,\ldots,x_n)\)，且 \(I\subseteq\mathfrak n\) 为理想。令
	\[
		R=(P/I)_{\mathfrak n/I},\qquad
		\mathfrak m=(\mathfrak n/I)R.
	\]
	则映射 \(X_i\mapsto[x_i]\in\mathfrak m/\mathfrak m^2\) 给出分次 \(k\) 代数同构
	\[
		\operatorname{gr}_{\mathfrak m}R
		\cong k[X_1,\ldots,X_n]/\operatorname{in}_0(I).
	\]
\end{proposition}
\begin{proof}
	先看满射性。极大理想 \(\mathfrak m\) 由 \(x_i\) 的像生成，所以 \(\mathfrak m^d\) 由总次数为 \(d\) 的单项式生成。模去 \(\mathfrak m^{d+1}\) 后，它们的系数只需取在 \(R/\mathfrak m=k\) 中。因此这些单项式张成第 \(d\) 层，所述分次映射满射。

	若 \(f\in I\) 的最低次数为 \(d\)，则 \(f=f^\ast+h\)、\(h\in\mathfrak n^{d+1}\)。在 \(R\) 中有 \(f=0\)，所以 \(f^\ast\) 的像属于 \(\mathfrak m^{d+1}\)，在第 \(d\) 层中为零。因此 \(\operatorname{in}_0(I)\) 包含于核。

	核是齐次理想，反向只须考察非零齐次多项式 \(F\)，设其次数为 \(d\)。若 \(F\) 在第 \(d\) 层的像为零，则
	\[
		F\in IP_{\mathfrak n}+\mathfrak n^{d+1}P_{\mathfrak n}.
	\]
	清去局部分母，可取 \(s\in P\)、\(s(0)\ne0\)，使
	\[
		sF=g+h,\qquad g\in I,\quad h\in\mathfrak n^{d+1}.
	\]
	由于 \(F\) 齐次，\(g=sF-h\) 的最低次项恰为 \(s(0)F\)。故 \(F=s(0)^{-1}g^\ast\in\operatorname{in}_0(I)\)，证明核等于该理想。
\end{proof}

对一个主理想 \(I=(f)\)，有
\[
	\operatorname{in}_0((f))=(f^\ast).
\]
这是因为多项式环为整环，非零乘积满足
\[
	\operatorname{ord}_0(hf)=\operatorname{ord}_0(h)+\operatorname{ord}_0(f),
	\qquad (hf)^\ast=h^\ast f^\ast.
\]
多个生成元之间则可能发生最低次项的抵消。

\begin{example}
\label{g5:exm:initial-generators-fail}
	设 \(k\) 为域，在 \(k[x,y]\) 中取
	\[
		I=(x,x+y^2)=(x,y^2).
	\]
	给定两个生成元的初始形式都是 \(X\)，它们只生成 \((X)\)。但两生成元相减得到 \(y^2\in I\)，所以 \(Y^2\in\operatorname{in}_0(I)\)。事实上，\(I=(x,y^2)\) 已是齐次理想，每个元素的齐次部分仍在 \(I\) 中，故
	\[
		\operatorname{in}_0(I)=(X,Y^2).
	\]
	只取给定生成元的初始形式，会把一个长度为二的加厚点误算成一整条直线。
\end{example}

\subsection{切锥与切空间}
\label{g5:subsec:0504:02}

最低次初始理想描述了一个齐次方程组。为了使它与坐标选择无关，应当把伴随分次环本身作为定义；幂零元也保留在这个环中。第一章对平面曲线所说的切锥，正是这一构造的主理想情形。

\begin{definition}
\label{g5:def:tangent-cone}
	设 \((R,\mathfrak m,\kappa)\) 为非零 Noether 局部环。以分次 \(\kappa\) 代数 \(\operatorname{gr}_{\mathfrak m}R\) 为坐标环的仿射锥，称为 \(R\) 在其闭点处的切锥，记为
	\[
		\mathrm C(R)=\operatorname{Spec}\bigl(\operatorname{gr}_{\mathfrak m}R\bigr).
	\]
	若 \(R=\mathcal O_{X,p}\)，也记为 \(\mathrm C_pX\)。这里的切锥包括指定坐标环的结构；只考察其约化支集时，会明确说明取根式。
\end{definition}
\symidx{\(\mathrm C_pX\)：点 \(p\) 处的切锥}

这个定义中的“锥”可由分次看出：若 \(\lambda\in\kappa\)，在次数 \(d\) 部分乘以 \(\lambda^d\) 与环乘法相容，给出标量伸缩。取局部环同构会把各次极大理想幂及其商一起对应，所以切锥不依赖仿射坐标的选择。方程式解释与平面曲线的切线方向可对照\cite[第2章，第1.5节，第94--95页]{Shafarevich2013Basic1}。

\ProseParagraphBreak

记 \(V=\mathfrak m/\mathfrak m^2\)，这是有限维 \(\kappa\) 向量空间。它的\term[duichengdaishu]{对称代数}[symmetric algebra] \(\operatorname{Sym}_\kappa(V)\) 是由 \(V\) 自由生成的交换 \(\kappa\) 代数：任一 \(\kappa\) 线性映射 \(V\to D\)，其中 \(D\) 为交换 \(\kappa\) 代数，都唯一延伸为代数同态 \(\operatorname{Sym}_\kappa(V)\to D\)。具体构造时，选 \(V\) 的一组基就得到以这些基向量为一次变量的多项式环；换基按同一个线性代换延伸，因而所得代数及泛性质与选基无关。对偶向量空间
\[
	\mathrm T(R)=\operatorname{Hom}_\kappa(V,\kappa)
\]
所对应的仿射空间，以 \(\operatorname{Sym}_\kappa(V)\) 为坐标环。对 \(k\) 有理点，这就是\cref{g5:thm:rational-cotangent-space}建立的切空间；一般剩余域情形，此处采用局部环自身的余切空间及其对偶。

\begin{proposition}
\label{g5:prop:tangent-cone-in-tangent-space}
	设 \((R,\mathfrak m,\kappa)\) 为非零 Noether 局部环。次数一部分的恒等映射唯一诱导满同态
	\[
		\operatorname{Sym}_\kappa(\mathfrak m/\mathfrak m^2)
		\twoheadrightarrow \operatorname{gr}_{\mathfrak m}R.
	\]
	其核为齐次理想，且零次、一次部分均为零。因此切锥由切空间中的齐次方程给出，其中没有非零一次方程。
\end{proposition}
\begin{proof}
	把 \(V=\mathfrak m/\mathfrak m^2\) 放入伴随分次环的次数一部分，由对称代数的泛性质得到所述映射。每个 \(\mathfrak m^d\) 中的元素都是有限多个 \(d\) 个极大理想元素乘积之和；模去 \(\mathfrak m^{d+1}\) 后，便由 \(V\) 中元素的 \(d\) 次乘积生成。所以该映射在每个次数都满射。

	映射保持次数，因而核齐次；在次数零它是 \(\kappa\) 的恒等，在次数一它是 \(V\) 的恒等，所以这两个次数的核均为零。商环的素谱对应给出切锥在切空间中的闭嵌入及所述方程。
\end{proof}

这里的结论保留了坐标环。若将切锥的方程理想取根式，可能出现新的线性方程；尖点的二重切线就是如此。另一方面，若 \(\mathfrak q\) 为分次环所有正次数部分的直和，则
\[
	\mathfrak q/\mathfrak q^2\cong\mathfrak m/\mathfrak m^2,
\]
因为分次环由次数一生成，所有次数至少二的项恰组成 \(\mathfrak q^2\)。所以切锥顶点的一阶切空间仍然是原来的切空间。

\subsection{尖点、结点和加厚点}
\label{g5:subsec:0504:03}

下面的三个模型在闭点处都有二维切空间。它们的区别来自伴随分次环中的二次及更高次关系，而不只是点集所包含的方向。

\begin{example}
\label{g5:exm:cusp-tangent-cone}
	设 \(k\) 为域，令
	\[
		R=\bigl(k[x,y]/(y^2-x^3)\bigr)_{(x,y)}.
	\]
	最低次项为 \(y^2\)，故
	\[
		\operatorname{gr}_{\mathfrak m}R
		\cong k[X,Y]/(Y^2).
	\]
	其约化支集是直线 \(Y=0\)，但 \([Y]\ne0\)、\([Y]^2=0\)。在这条直线的泛点处，局部化得到
	\[
		k(X)[Y]/(Y^2),
	\]
	这是长度为二的局部环；“二重切线”指的正是保留了这一厚度。它不是两条不同直线。

	次数零部分为 \(k\)。对每个 \(d\geq1\)，次数 \(d\) 部分的一组基为
	\[
		X^d,\quad X^{d-1}Y,
	\]
	所以其维数恒为二。虽然原局部环是整环，伴随分次环已经出现非零幂零元。
\end{example}

\begin{example}
\label{g5:exm:node-tangent-cone}
	设 \(k\) 为特征不等于二的域，令
	\[
		R=\bigl(k[x,y]/(y^2-x^2-x^3)\bigr)_{(x,y)}.
	\]
	则
	\[
		\operatorname{gr}_{\mathfrak m}R
		\cong k[X,Y]/(Y^2-X^2)
		=k[X,Y]/\bigl((Y-X)(Y+X)\bigr).
	\]
	两个线性因子不相伴，故理想是两个素理想 \((Y-X)\)、\((Y+X)\) 的交：同时被这两个互素多项式整除的多项式，被它们的乘积整除。因此这个分次环约化，切锥是两条不同直线的并。

	关于首一的 \(Y\) 二次式作除法，每个次数 \(d\geq1\) 的齐次部分仍以 \(X^d,X^{d-1}Y\) 为基。它与尖点的各层维数相同，但环结构不同：一个约化，另一个不约化。特征为二时，这个最低次方程成为 \((Y-X)^2\)，不能再作两条不同直线的解释。
\end{example}

\begin{example}
\label{g5:exm:thick-point-tangent-cone}
	设 \(k\) 为域，取
	\[
		R=k[x,y]/(x,y)^2,\qquad \mathfrak m=(x,y).
	\]
	这个环已经是局部环，\(\mathfrak m^2=0\)，并以 \(1,x,y\) 为 \(k\) 基。其伴随分次环为
	\[
		\operatorname{gr}_{\mathfrak m}R
		=k\oplus(kX\oplus kY)
		\cong k[X,Y]/(X,Y)^2.
	\]
	切空间为二维平面；切锥的约化支集却只有原点。整个切锥环的长度为三，次数零、一次、二次以后的维数分别为 \(1,2,0\)。这三个组成因子与一条二重直线的正维支集是不同的信息。
\end{example}

\begin{table}[htbp]
	\centering
	\caption{三个二维切空间模型的切锥}
	\label{g5:tab:tangent-cone-models}
	\begin{tabular}{llll}
		\toprule
		模型 & 切锥坐标环 & 约化支集 & 各层维数 \(d=0,1,2,\ldots\) \\
		\midrule
		尖点 & \(k[X,Y]/(Y^2)\) & 一条直线 & \(1,2,2,2,\ldots\) \\
		结点（\(\operatorname{char}k\ne2\)） & \(k[X,Y]/(Y^2-X^2)\) & 两条直线 & \(1,2,2,2,\ldots\) \\
		加厚点 & \(k[X,Y]/(X,Y)^2\) & 一个点 & \(1,2,0,0,\ldots\) \\
		\bottomrule
	\end{tabular}
\end{table}

\cref{g5:tab:tangent-cone-models}也区分了三个比较层次：切空间只读取一次层，各层维数只计数齐次向量空间，切锥则还保留这些层之间的乘法。即使切锥相同，高阶项仍可能改变形式方程；要研究这些项，需要把各阶截断相容地放在一起。

\subsection{形式同构与有限阶截断的差异}
\label{g5:subsec:0504:04}

伴随分次保留各层及其乘法，但把一个元素属于哪一层的高阶修正分开了。要把这些修正也记录下来，可以同时考察
\[
	R/\mathfrak m,\quad R/\mathfrak m^2,\quad R/\mathfrak m^3,\quad\ldots
\]
及它们之间的取商映射。一个有限截断只保留有限多阶；整个相容系统则允许逐阶确定一个形式表达式。

\begin{definition}
\label{g5:def:formal-series-and-completion}
	设 \(k\) 为域。变量 \(x=(x_1,\ldots,x_n)\) 的\term[xingshimijishuhuan]{形式幂级数环}[formal power series ring] \(k[[x]]\) 由所有表达式
	\[
		F=\sum_{\alpha\in\mathbb N^n}c_\alpha x^\alpha,\qquad c_\alpha\in k
	\]
	组成。加法逐项进行，乘法中 \(x^\alpha\) 的系数为有限和
	\[
		\sum_{\beta+\gamma=\alpha}c_\beta d_\gamma.
	\]
	设 \((R,\mathfrak m)\) 为局部环。各商环及自然映射的\term[nijixian]{逆极限}[inverse limit]在此具体指环
	\[
		\widehat R
		=\varprojlim_{q\geq1}R/\mathfrak m^q
		=\left\{(a_q)_{q\geq1}:
		\begin{array}{l}
			a_q\in R/\mathfrak m^q,\\
			a_{q+1}\bmod\mathfrak m^q=a_q
		\end{array}
		\right\},
	\]
	其运算逐个分量进行。它称为 \(R\) 的极大理想进\term[wanbeihua]{完备化}[completion]，自然同态为 \(r\mapsto(r\bmod\mathfrak m^q)_q\)。
\end{definition}
\symidx{\(k[[x_1,\ldots,x_n]]\)：形式幂级数环}
\symidx{\(\widehat R=\varprojlim_q R/\mathfrak m^q\)：局部环的完备化}

形式级数的每个齐次部分都是有限项多项式，非零级数仍以最低非零齐次次数定义 \(\operatorname{ord}_0\)。若记 \(\mathfrak n=(x_1,\ldots,x_n)\subseteq k[[x]]\)，则 \(\mathfrak n^q\) 恰由最低总次数至少为 \(q\) 的级数组成：将每个这样的单项式分配给它的一个 \(q\) 次单项式因子，再按有限多个可能的因子分组即可。因此一个相容截断序列唯一确定全部系数，给出
\[
	k[[x]]\cong\varprojlim_q k[x]/(x_1,\ldots,x_n)^q.
\]
这里的逐阶收敛指每个固定次数的系数最终稳定，不涉及数值意义下的级数收敛。若 \(F\) 的常数项为 \(c\ne0\)，写 \(F=c(1-H)\)、\(H\in\mathfrak n\)，则
\[
	F^{-1}=c^{-1}\sum_{j\geq0}H^j
\]
按上述规则有意义；反之常数项为零的元素不可能可逆。所以 \(k[[x]]\) 是极大理想为 \(\mathfrak n\) 的局部环。定义及逐阶展开可对照\cite[第2章，第2.2节，定义2.4，第101页]{Shafarevich2013Basic1}；逆极限定义可对照\cite[第8节，第55页]{Matsumura1989}。

\ProseParagraphBreak

在 \(\widehat R\) 上，把
\[
	F^q\widehat R=\ker(\widehat R\longrightarrow R/\mathfrak m^q)
\]
及其平移作为开邻域，便得到由截断规定的拓扑。对局部 \(k\) 代数 \(R,S\)，若存在两侧都连续的 \(k\) 代数同构 \(\widehat R\cong\widehat S\)，就称它们\term[xingshitonggou]{形式同构}[formally isomorphic]。当它们是点的局部环时，也称这两个点的形式局部类型相同。下面先把平面方程的完备化写成可计算的形式商环。

\begin{proposition}
\label{g5:prop:principal-formal-completion}
	设 \(k\) 为域，\(P=k[x_1,\ldots,x_n]\)，\(\mathfrak n=(x_1,\ldots,x_n)\)，且 \(0\ne f\in\mathfrak n\)。令
	\[
		R=(P/(f))_{\mathfrak n/(f)}.
	\]
	则逐阶取系数给出同构
	\[
		\widehat R\cong k[[x_1,\ldots,x_n]]/(f).
	\]
	右侧采用极大理想幂的拓扑时，这也是拓扑环同构。
\end{proposition}
\begin{proof}
	记 \(S=k[[x_1,\ldots,x_n]]\)、\(\mathfrak n_S=(x_1,\ldots,x_n)S\)。常数项非零的多项式在模 \(\mathfrak n^q\) 后可用有限几何级数求逆，故局部化不再改变这一截断：
	\[
		R/\mathfrak m^q\cong P/\bigl((f)+\mathfrak n^q\bigr).
	\]
	于是有自然同态
	\[
		S\longrightarrow\varprojlim_q P/\bigl((f)+\mathfrak n^q\bigr).
	\]
	我们直接证明它满射且核为 \((f)\)。

	给定一个相容序列，归纳选择代表多项式 \(p_q\)。若 \(p_q\) 已选好，任取下一项的代表 \(b\)，相容性给出
	\[
		b-p_q=fh+u,\qquad h\in P,\quad u\in\mathfrak n^q.
	\]
	取 \(p_{q+1}=b-fh\)，它仍代表下一项，而且 \(p_{q+1}-p_q\in\mathfrak n^q\)。因此各固定次数的系数稳定，确定一个 \(F\in S\)，其像就是所给序列，证明满射。

	再设 \(F\) 落在核中。对每个 \(q\)，可取 \(g_q\in P\) 使
	\[
		F-fg_q\in\mathfrak n_S^q.
	\]
	令 \(d=\operatorname{ord}_0(f)\)。在形式幂级数环中，两个非零级数的最低次齐次部分之积仍非零，所以
	\[
		\operatorname{ord}_0(fG)=d+\operatorname{ord}_0(G)
	\]
	对非零 \(G\in S\) 成立。因而，当 \(r\geq q>d\) 时，由
	\(f(g_r-g_q)\in\mathfrak n_S^q\)
	可得 \(g_r-g_q\in\mathfrak n_S^{q-d}\)。这使 \(g_q\) 的各固定次数系数最终稳定，确定级数 \(G\in S\)。逐阶取系数得到 \(F=fG\)。反向包含显然成立，因此核恰为 \((f)\)。

	所构造的同构把投影到 \(R/\mathfrak m^q\) 的核对应到
	\((\mathfrak n_S^q+(f))/(f)\)，即右侧极大理想的 \(q\) 次幂。因此它及其逆都连续。
\end{proof}

这项证明只使用代表元的逐阶调整和主理想的闭合性计算。在这里，后一个计算实际上只用到乘以非零 \(f\) 会使最低次数恰好增加固定的 \(d\)。以下形式坐标变换也按同样的逐阶规则理解：把常数项为零的级数代入另一个级数，每个指定次数仍只涉及有限项。

\begin{example}
\label{g5:exm:formal-node}
	设 \(k\) 为特征不等于二的域。结点
	\[
		y^2-x^2(1+x)=0
	\]
	在原点的形式局部环同构于 \(k[[u,v]]/(uv)\)。

	为构造这个同构，先求常数项为一的级数
	\[
		h(x)=1+\sum_{n\geq1}a_nx^n,\qquad h(x)^2=1+x.
	\]
	比较系数得到
	\[
		2a_1=1,\qquad
		2a_n=-\sum_{i=1}^{n-1}a_i a_{n-i}\quad(n\geq2).
	\]
	因为只需除以 \(2\)，这个递推在所有特征不为二的域上都有唯一解。令
	\[
		u=y-xh(x),\qquad v=y+xh(x).
	\]
	则
	\[
		uv=y^2-x^2h(x)^2=y^2-x^2(1+x).
	\]

	还须验证它确实是可逆的形式坐标变换。记 \(q(x)=xh(x)=x+\) 高次项。逐阶构造
	\(\phi(z)=z+\sum_{n\geq2}b_nz^n\)
	使 \(q(\phi(z))=z\)：当 \(b_2,\ldots,b_{n-1}\) 已知时，次数 \(n\) 的系数是 \(b_n\) 加上已知项，因而唯一确定 \(b_n\)。又 \(q(U)-q(V)\) 可写成 \((U-V)\) 乘以常数项为一的级数，所以 \(q(U)=q(V)\) 必有 \(U=V\)。将 \(U=\phi(q(x))\)、\(V=x\) 代入，得到 \(\phi(q(x))=x\)。因此逆变换为
	\[
		y=\frac{u+v}{2},\qquad
		x=\phi\left(\frac{v-u}{2}\right).
	\]
	这给出形式幂级数环的互逆同构，并把方程理想对应起来，故由上个命题得到
	\[
		\widehat{\bigl(k[x,y]/(y^2-x^2-x^3)\bigr)_{(x,y)}}
		\cong k[[u,v]]/(uv).
	\]
	右侧有两个极小素理想 \((u)\)、\((v)\)，因为任何包含 \(uv\) 的素理想都包含其中一个因子，而相应商环均为一元形式幂级数整环。这就是结点的两个形式分支。
\end{example}

第四章已经证明原结点的局部环是整环；形式局部环却有两个分支。形式展开改变了可使用的函数，所以这里的分支分解与原局部环中的素理想分解不同。此例可对照\cite[第I章，第5节，例5.6.3，第34--35页]{Hartshorne1977}；上述递推把平方根与坐标逆变换都明确写出。

\begin{example}
\label{g5:exm:finite-truncation-not-formal-type}
	设 \(k\) 为域，给定任意整数 \(N\geq1\)，取
	\[
		R=k[x]_{(x)},\qquad S=k[x]/(x^{N+1}).
	\]
	对每个 \(1\leq q\leq N+1\)，两侧的截断都为 \(k[x]/(x^q)\)，这些有限个同构还与取商映射相容。但
	\[
		\widehat R=k[[x]],\qquad \widehat S=S,
	\]
	因为前者由全部系数构成，后者的截断系统从 \(q=N+1\) 起已稳定。第一个环为整环，第二个环中 \(x\ne0\) 且 \(x^{N+1}=0\)，所以它们不形式同构。已知一个指定的有限阶截断，不能对所有局部环决定其形式局部类型。
\end{example}

若已知的是全部截断及其相容关系，则情况不同。关键条件是同构能够在各阶之间组成一个系统，而不仅是每一阶分别存在某个同构。

\begin{proposition}
\label{g5:prop:compatible-truncation-isomorphisms}
	设 \((R,\mathfrak m)\)、\((S,\mathfrak n)\) 为局部 \(k\) 代数。假设对每个 \(q\geq1\)，都有 \(k\) 代数同构
	\[
		\varphi_q:R/\mathfrak m^q\xrightarrow{\ \sim\ }S/\mathfrak n^q,
	\]
	并且以下方块交换：
	\[
		\begin{tikzcd}
			R/\mathfrak m^{q+1} \arrow[r,"\varphi_{q+1}"] \arrow[d]
			& S/\mathfrak n^{q+1} \arrow[d] \\
			R/\mathfrak m^q \arrow[r,"\varphi_q"]
			& S/\mathfrak n^q.
		\end{tikzcd}
	\]
	则它们唯一诱导一个与这些投影相容的拓扑 \(k\) 代数同构
	\(\widehat R\cong\widehat S\)。
\end{proposition}
\begin{proof}
	对 \((a_q)_q\in\widehat R\)，定义其像为 \((\varphi_q(a_q))_q\)。方块交换保证这些分量相容，故像属于 \(\widehat S\)。逐分量运算说明这是 \(k\) 代数同态。各 \(\varphi_q^{-1}\) 也组成相容系统，按同样规则给出逆同态。两侧投影核逐一对应，所以两个方向都连续。任一与投影相容的同态在第 \(q\) 个分量上都必须等于 \(\varphi_q\)，因而唯一。
\end{proof}

切空间、切锥、有限阶截断和形式局部环由此形成不同层次的数据。切空间只有一次信息；切锥保留各次初始形式及其乘法；有限截断保留有限阶内的高阶修正；完备化则把所有相容截断同时保留下来。它们之间的比较需要检验具体环结构，不能只比较各层的维数。

% 纲要条目排为编号小节标题，并进入目录与书签；说明段仍保留于纲要注释中。


\begin{chaptersummary}
	微分模把以任意模为值域的导子统一表示为线性映射。多项式环的坐标微分自由生成，局部化给出分式求导规则；商环加入方程的微分关系。余法正合列描述这些关系，传递序列则把底环中允许变化的元素逐步设为常数。两种序列的左端都可能有核，微分模的基变换同构也不保证余法模在任意基变换下保持不变。

	\ProseParagraphBreak

	在 \(k\) 有理点处，双数点的乘法条件正是 乘法公式，切空间由 \(\mathfrak m/\mathfrak m^2\) 的对偶给出。Jacobian 的秩计算这一空间的维数；与同一点的局部维数比较后，才能判定正则性。一般剩余域还具有自身相对于 \(k\) 的微分，不可分扩张会使这种差别显现出来。

	\ProseParagraphBreak

	切锥保留极大理想各次幂之间的乘法，而切空间只保留第一层。计算切锥需要整个理想的最低次初始形式；任取一组方程后分别取最低次项，可能丢失相消产生的新关系。全部相容的有限阶截断形成形式邻域；一个有限阶截断可以相同，却不足以确定形式方程。尖点、结点和加厚点显示了这些不同层次各自保存的几何信息。
\end{chaptersummary}

\BookExercises

\begin{exercise}\label{g5:ex:laurent-derivations}
	设 \(k\) 为域，\(B=k[t,t^{-1}]\)。确定所有 \(k\) 导子 \(B\to B\)，并计算由 \(E(t)=t\) 确定的导子 \(E\) 的核，分别讨论特征零和特征 \(p>0\)。
\end{exercise}
\begin{solution}
	任意导子由 \(D(t)=h\in B\) 唯一确定，且
	\(D(t^{-1})=-t^{-2}h\)。反复使用乘法公式，对全部 \(n\in\mathbb Z\) 都有
	\(D(t^n)=nt^{n-1}h\)。反过来，这个公式逐项定义的映射满足导子条件，故
	\(\operatorname{Der}_k(B,B)=B\,\partial/\partial t\)。
	若 \(f=\sum_n a_nt^n\)，则 \(E(f)=\sum_n na_nt^n\)。Laurent 单项式线性无关，所以特征零时核为 \(k\)，特征 \(p\) 时核为 \(k[t^p,t^{-p}]\)。这个结论不要求 \(k\) 完美：核的系数可以是任意 \(k\) 中元素。
\end{solution}

\begin{exercise}\label{g5:ex:truncated-polynomial-differentials}
	设 \(k\) 为域，\(n\geq2\)，\(B=k[t]/(t^n)\)。计算 \(\Omega_{B/k}\)、\(\operatorname{Der}_k(B,B)\) 及它们的 \(k\) 维数，并比较 \(\Omega_{B/k}\otimes_B k\) 的维数。
\end{exercise}
\begin{solution}
	由方程的微分关系，
	\[
		\Omega_{B/k}=B\,\mathrm{d}t/(nt^{n-1}\mathrm{d}t).
	\]
	导子由 \(h=D(t)\) 确定，而下降条件为 \(nt^{n-1}h=0\)。若 \(n\) 在 \(k\) 中非零，微分模为 \((B/(t^{n-1}))\,\mathrm{d}t\)，导子的 \(h\) 恰为 \((t)\) 中元素，两者维数都是 \(n-1\)。若 \(\operatorname{char}k\mid n\)，微分模自由于 \(\mathrm{d}t\)，任意 \(h\in B\) 都允许，两者维数都是 \(n\)。在两种情形中，模去 \((t)\) 后微分关系均为零，故微分模的纤维恒为一维。纤维维数不能恢复整个微分模的长度。
\end{solution}

\begin{exercise}\label{g5:ex:square-zero-differentials}
	设 \(k\) 为特征不为二的域，\(V\) 为 \(r\) 维 \(k\) 向量空间。在 \(B=k\oplus V\) 上规定
	\((a,u)(b,v)=(ab,av+bu)\)。计算 \(\dim_k\Omega_{B/k}\)，并证明其在唯一极大理想 \(V\) 处的纤维维数为 \(r\)。
\end{exercise}
\begin{solution}
	选 \(V\) 的基 \(x_1,\ldots,x_r\)，则
	\(B=k[x_1,\ldots,x_r]/(x_ix_j:1\leq i,j\leq r)\)。
	微分模由 \(\mathrm{d}x_j\) 生成，关系为
	\(x_i\,\mathrm{d}x_j+x_j\,\mathrm{d}x_i=0\)，特别有 \(x_i\,\mathrm{d}x_i=0\)。
	因此下列元素张成它作为 \(k\) 向量空间：
	\[
		\mathrm{d}x_i\ (1\leq i\leq r),\qquad
		x_i\,\mathrm{d}x_j\ (i<j).
	\]

	为验证没有额外关系，在 \(N=V\oplus\bigwedge^2V\) 上定义
	\[
		(a,u)\cdot(v,w)=(av,aw+u\wedge v).
	\]
	两次作用的复合符合 \(B\) 的乘法，故这是 \(B\) 模。映射
	\(D(a,u)=(u,0)\) 为 \(k\) 导子，因为乘法公式第二分量中的 \(u\wedge v+v\wedge u\) 为零。它使上述生成元分别映到 \(x_i\) 和 \(x_i\wedge x_j\)，这些元素为 \(N\) 的基，故上述张成组也线性无关。因此
	\[
		\dim_k\Omega_{B/k}=r+\binom r2=\frac{r(r+1)}2.
	\]
	\(V\Omega_{B/k}\) 正是第二类元素张成的空间，除去它后只剩 \(r\) 个坐标微分。\(B\) 的每个 \(a+u\)、\(a\ne0\) 都有逆 \(a^{-1}-a^{-2}u\)，所以 \(V\) 确为唯一极大理想。
\end{solution}

\begin{exercise}\label{g5:ex:product-differentials}
	设 \(A\) 为交换环，\(B,C\) 为交换 \(A\) 代数，\(R=B\times C\)。先证明任意导子都杀死幂等元，再证明存在自然的 \(R\) 模同构
	\[
		\Omega_{R/A}\cong\Omega_{B/A}\times\Omega_{C/A},
	\]
	其中 \(R\) 在右边按两个投影分别作用。
\end{exercise}
\begin{solution}
	若 \(e^2=e\)，则 \(D(e)=2eD(e)\)。再乘 \(e\) 得 \(eD(e)=2eD(e)\)，消去后有 \(eD(e)=0\)，从而 \(D(e)=0\)。这一步对任意特征都成立。

	在 \(R\) 中取 \(e=(1,0)\)、\(f=(0,1)\)。任意 \(R\) 模 \(M\) 都分解为 \(eM\oplus fM\)：由 \(e+f=1\) 得到和，由 \(ef=0\) 得到交为零。一个 \(A\) 导子 \(D:R\to M\) 满足
	\[
		D(b,0)=eD(b,0),\qquad D(0,c)=fD(0,c).
	\]
	这是对 \((b,0)=e(b,0)\) 和 \((0,c)=f(0,c)\) 使用乘法公式的结果。因此它的两个限制分别取值于 \(eM,fM\)，并在 \(A\) 的像上为零：例如 \((a_B,0)=e(a_B,a_C)\)，而右边两因子的导数都为零。故两个限制分别为 \(B\) 与 \(C\) 的 \(A\) 导子。

	反过来，给定这两个导子，按 \(D(b,c)=D_B(b)+D_C(c)\) 定义就得到 \(R\) 上的导子，混合乘积为零与 \(ef=0\) 相容。分别应用两个微分模的泛性质，得到
	\(\Omega_{B/A}\times\Omega_{C/A}\) 连同 \((b,c)\mapsto(\mathrm{d}b,\mathrm{d}c)\) 满足 \(\Omega_{R/A}\) 的泛性质，故所述同构成立。
\end{solution}

\begin{exercise}\label{g5:ex:node-family-differentials}
	设 \(k\) 为域，\(A=k[s]\)、\(B=k[x,y]\)，结构同态由 \(s\mapsto xy\) 给出。写出 \(k\to A\to B\) 的传递序列，证明其左端单射；再计算纤维 \(B_a=B/(xy-a)\)、\(a\in k\)，相对于 \(k\) 的微分模。说明 \(a=0\) 时该模在原点不自由。
\end{exercise}
\begin{solution}
	传递序列为
	\[
		B\,\mathrm{d}s\longrightarrow
		B\,\mathrm{d}x\oplus B\,\mathrm{d}y
		\longrightarrow\Omega_{B/A}\longrightarrow0,
		\qquad \mathrm{d}s\longmapsto y\,\mathrm{d}x+x\,\mathrm{d}y.
	\]
	\(B\) 是整环且 \(x,y\ne0\)，故第一张映射单射。
	由基变换公式或直接用商环呈示，
	\[
		\Omega_{B_a/k}
		=\frac{B_a\,\mathrm{d}x\oplus B_a\,\mathrm{d}y}
		{B_a(y\,\mathrm{d}x+x\,\mathrm{d}y)}.
	\]
	当 \(a\ne0\) 时 \(x\) 为单位，\(y=a/x\)，于是该模自由于 \(\mathrm{d}x\)。
	当 \(a=0\) 时，在原点剩余域上的纤维为二维；若该局部模自由，其秩必为二。再在素理想 \((x)\subset k[x,y]/(xy)\) 处局部化，\(y\) 可逆，\(x=0\)，局部环为 \(k(y)\)，而微分模为一维 \(k(y)\,\mathrm{d}y\)。自由秩二会在此仍为二，矛盾。因此原点处不自由。
\end{solution}

\begin{exercise}\label{g5:ex:quadratic-cone-tangent-map}
	设 \(k\) 为特征不为二的域，\(X\) 为 \(xz-y^2=0\) 定义的仿射曲面。对映射
	\[
		F:\mathbb A_k^2\longrightarrow X,\qquad
		(u,v)\longmapsto(u^2,uv,v^2),
	\]
	计算各 \(k\) 有理点的切映射。比较原点与非原点的切空间，并证明 \(X\) 的非原点 \(k\) 有理点正则，而原点不正则。
\end{exercise}
\begin{solution}
	先代入方程验证 \(u^2v^2-(uv)^2=0\)。切映射为
	\[
		(a,b)\longmapsto(2ua,va+ub,2vb).
	\]
	在 \((u,v)=(0,0)\) 时它为零；在非原点处为单射：若 \(u\ne0\)，第一式迫使 \(a=0\)，第二式迫使 \(b=0\)；\(v\ne0\) 时同理。
	目标点 \((x,y,z)\) 的切方程为 \(z\xi-2y\eta+x\zeta=0\)，所以原点切空间三维，非原点切空间二维。

	在 \(X\) 的任一非原点 \(k\) 有理点，\(x,z\) 至少有一个非零。例如 \(x\ne0\) 时坐标环局部化为 \(k[x,x^{-1},y]\)，因为 \(z=y^2/x\)。该有理点的局部极大理想由 \(x-x_0,y-y_0\) 生成，且有长度二的素链，故局部维数恰为二，因而正则。原点处局部维数也为二：超曲面维数给出上界二，而
	\((0)\subset(x,y)\subset(x,y,z)\) 给出下界。这里坐标环是整环，可由 \(xz-y^2\) 在 \(k[x,y][z]\) 中为本原、在 \(k(x,y)[z]\) 中为一次多项式的不可约性验证。原点嵌入维数三大于局部维数二，故不正则。
\end{solution}

\begin{exercise}\label{g5:ex:nonrational-separable-cotangent}
	取 \(R=\mathbb R[t]_{(t^2+1)}\)，\(\mathfrak m=(t^2+1)R\)，剩余域 \(\kappa\cong\mathbb C\)，其中 \(t\) 的剩余类对应 \(i\)。计算
	\[
		\mathfrak m/\mathfrak m^2
		\longrightarrow\Omega_{R/\mathbb R}\otimes_R\kappa
		\longrightarrow\Omega_{\kappa/\mathbb R}\longrightarrow0.
	\]
\end{exercise}
\begin{solution}
	\(R\) 是以 \(t^2+1\) 为统一化元的 DVR，故第一项为一维 \(\kappa\) 空间，以 \(t^2+1\) 的类为基。局部化公式使中间项为 \(\mathbb C\,\mathrm{d}t\)，第一张映射把该基送到 \(2i\,\mathrm{d}t\)，因而是同构。最后一项为零，也可直接由最小多项式 \(T^2+1\) 的导数 \(2i\ne0\) 得到。此处剩余域不是基域，但其可分性使剩余域本身没有相对微分。
\end{solution}

\begin{exercise}\label{g5:ex:inseparable-plane-cotangent}
	设 \(k=\mathbb F_p(s)\)，\(A=k[x,y]\)，\(\mathfrak p=(x^p-s,y)\)，\(R=A_{\mathfrak p}\)，\(\mathfrak m=\mathfrak pR\)，\(\kappa=R/\mathfrak m\)。证明 \(R\) 为二维正则局部环，并计算余切正合尾。说明即使 \(\mathfrak m/\mathfrak m^2\) 与 \(\Omega_{A/k}\otimes_A\kappa\) 维数相同，它们之间的自然映射仍可能不是同构。
\end{exercise}
\begin{solution}
	记 \(f=x^p-s\)。由\cref{g5:exm:inseparable-regular-point}的域扩张计算，\(f\) 在 \(k[x]\) 中不可约，所以 \(\mathfrak p\) 为极大理想，\(\kappa=k[\alpha]\)、\(\alpha^p=s\)。局部化保留严格素链
	\[
		(0)\subsetneq(y)\subsetneq(f,y).
	\]
	因而 \(\dim R\geq2\)，而多项式环维数给出 \(\dim R\leq2\)。极大理想由 \(f,y\) 生成，所以嵌入维数至多为二；维数下界使它恰为二，\(R\) 正则，且 \([f],[y]\) 构成 \(\mathfrak m/\mathfrak m^2\) 的一组基。

	余切正合尾为
	\[
		\kappa\cdot[f]\oplus\kappa\cdot[y]
		\xrightarrow{\alpha_0}
		\kappa\,\mathrm{d}x\oplus\kappa\,\mathrm{d}y
		\longrightarrow\kappa\,\mathrm{d}\alpha\longrightarrow0,
	\]
	其中 \(\alpha_0([f])=0\)、\(\alpha_0([y])=\mathrm{d}y\)，后一张映射把 \(\mathrm{d}x\) 送到 \(\mathrm{d}\alpha\)，把 \(\mathrm{d}y\) 送到零。因此前两个空间虽然均为二维，连接它们的自然映射却只有秩一。只有维数相同，并不能取代同构的验证；剩余域自身的非零微分正好记录这里的余核。
\end{solution}

\begin{exercise}\label{g5:ex:tangent-vector-higher-lift}
	设 \(k\) 为域。在原点处，双数点 \(x\mapsto a\varepsilon,y\mapsto b\varepsilon\) 对 \(xy=0\) 与 \(y^2-x^3=0\) 都成立。分别求它们延拓到 \(k[\varepsilon]/(\varepsilon^3)\) 的条件；对尖点再求延拓到 \(k[\varepsilon]/(\varepsilon^4)\) 的条件。
\end{exercise}
\begin{solution}
	写 \(x=a\varepsilon+c\varepsilon^2\)、\(y=b\varepsilon+d\varepsilon^2\)。对结点，模 \(\varepsilon^3\) 的方程为 \(ab\varepsilon^2=0\)，故条件是 \(ab=0\)，满足时任取 \(c,d\) 都可。对尖点，方程模 \(\varepsilon^3\) 为 \(b^2\varepsilon^2=0\)，因 \(k\) 为域，这等价于 \(b=0\)。

	模 \(\varepsilon^4\) 时，允许在 \(x,y\) 中加入三次项；已有 \(b=0\) 后，\(y^2\) 的阶至少为四，而 \(x^3\) 的三次项为 \(a^3\varepsilon^3\)。因此必须 \(a=0\)，且此时任取二次、三次系数都满足方程。于是尖点的每个非零切向量最终都遇到延拓障碍；参数化 \(x=t^2,y=t^3\) 的一阶导数为零，恰与此相容。
\end{solution}

\begin{exercise}\label{g5:ex:initial-ideal-extra-relation}
	设 \(k\) 为域，\(I=(x^2+y^3,xy)\subset k[x,y]\)，考察原点处的局部环。计算其最低次初始理想、切锥和各次齐次分量的维数，并指出给定两个生成元的最低次项漏掉了什么关系。
\end{exercise}
\begin{solution}
	消去得到 \(y(x^2+y^3)-x(xy)=y^4\in I\)，故初始理想包含
	\(J=(x^2,xy,y^4)\)。在商环中 \(x^2=-y^3\)、\(xy=0\)、\(y^4=0\)，所以 \(1,x,y,y^2,y^3\) 张成它。它们也线性无关：在五维向量空间上规定上述乘法、\(1\) 为单位、\(y^iy^j=y^{i+j}\)（指数至少四时为零），就得到一个结合代数，且由 \(x,y\) 生成并满足原方程。结合性只须检查含两个或更多 \(x\) 的乘积，均由 \(xy=0,x^2=-y^3\) 化为一致的零。由此构成与商环互逆的映射。

	\(x,y\) 均幂零，所以该环本已局部。其极大理想幂为
	\[
		\mathfrak m=(x,y),\quad
		\mathfrak m^2=(y^2,y^3),\quad
		\mathfrak m^3=(y^3),\quad
		\mathfrak m^4=0.
	\]
	伴随分次各层维数为 \(1,2,1,1\)。\(k[x,y]/J\) 也有这些维数，而它自然满射到伴随分次，故该满射逐次为同构，初始理想恰为 \(J\)。切锥为这个分次环所给的加厚点；遗漏的 \(y^4\) 来自最低次项相消后的新方程。
\end{solution}

\begin{exercise}\label{g5:ex:same-cone-different-local-rings}
	设 \(k\) 为域，
	\[
		R=k[x,y]/(x^2,xy,y^4),\qquad
		S=k[x,y]/(x^2-y^3,xy,y^4).
	\]
	证明它们都是长度五的局部 \(k\) 代数，有同构的伴随分次环，却不是同构的局部环。可比较被极大理想整体零化的子空间。
\end{exercise}
\begin{solution}
	两环都以 \(1,x,y,y^2,y^3\) 为基，验证方法与上一题的显式五维代数相同；全部非常数基向量均幂零，故各自唯一极大理想为 \((x,y)\)。两环的极大理想幂均为
	\(\mathfrak m^2=(y^2,y^3)\)、\(\mathfrak m^3=(y^3)\)、\(\mathfrak m^4=0\)。在 \(S\) 的伴随分次中，\(x^2=y^3\) 在二次层变成 \(x^2=0\)，并仍有 \(xy=0,y^4=0\)。维数比较给出
	\(\operatorname{gr}S\cong R\cong\operatorname{gr}R\)，右边按 \(x,y\) 的标准次数分次。

	写任意元素为 \(a+bx+cy+dy^2+ey^3\)。在 \(R\) 中同时乘 \(x,y\) 得零，等价于 \(a=c=d=0\)，所以被极大理想零化的子空间为 \(kx\oplus ky^3\)，维数二。在 \(S\) 中额外有 \(x\cdot bx=by^3\)，故还必须 \(b=0\)，相应子空间仅为 \(ky^3\)，维数一。任何局部环同构都保持极大理想及其零化子，而且诱导剩余域同构，故保持这一剩余域维数；两环不可能同构。它们已是有限阶的局部环，全部形式信息仍不能从共同的切锥恢复。
\end{solution}

\begin{exercise}\label{g5:ex:formal-one-variable-automorphism}
	设 \(k\) 为域，\(f(t)=a_1t+a_2t^2+\cdots\in t k[[t]]\)。证明代入 \(t\mapsto f(t)\) 给出保持 \((t)\) 进拓扑的 \(k\) 代数自同构，当且仅当 \(a_1\ne0\)。说明这些自同构如何对应于所有截断 \(k[t]/(t^{n+1})\) 上相容的局部自同构。
\end{exercise}
\begin{solution}
	\(f\) 的常数项为零，故模 \(t^{n+1}\) 计算任意复合只涉及有限多个系数，因此代入良定义且连续。若 \(a_1=0\)，它在 \((t)/(t^2)\) 上为零，不可能可逆。

	若 \(a_1\ne0\)，逐次构造 \(g(t)=b_1t+b_2t^2+\cdots\) 使 \(f(g(t))=t\)。一次方程给出 \(b_1=a_1^{-1}\)；已知 \(b_1,\ldots,b_{n-1}\) 后，\(n\) 次系数中未知量只以 \(a_1b_n\) 出现，其余由低阶系数决定，所以 \(b_n\) 存在且唯一。相同方法构造 \(h\) 使 \(h(f(t))=t\)。结合复合的结合性，有
	\(h=h\circ(f\circ g)=(h\circ f)\circ g=g\)，故 \(g\) 是两侧逆。

	一个连续自同构显然诱导相容的截断自同构。反过来，相容的截断映射把 \(t\) 送到一列相容的多项式类，唯一给出常数项为零的级数 \(f\)，且第一个非平凡截断可逆迫使 \(a_1\ne0\)。截断中的每个元素都是 \(t\) 的多项式，因此该级数代入恢复原来的全部截断映射。相容性正是各阶系数能同时选定的条件。
\end{solution}
